\documentclass[11pt,letterpaper]{article}
\usepackage[T1]{fontenc}
\usepackage[utf8]{inputenc}
\usepackage{lmodern,microtype}
\usepackage[margin=.92in,headheight=15pt]{geometry}
\usepackage{amsmath,amssymb,amsthm,mathtools}
\usepackage{booktabs,longtable,tabularx,array,enumitem}
\usepackage{xcolor,fancyhdr,listings,xurl,needspace}
\usepackage[colorlinks=true,linkcolor=blue!35!black,citecolor=blue!35!black,urlcolor=blue!45!black]{hyperref}
\usepackage{bookmark}
\definecolor{ink}{RGB}{25,46,67}
\definecolor{pale}{RGB}{246,248,250}
\setlength{\parindent}{0pt}
\setlength{\parskip}{.35em}
\setlength{\emergencystretch}{3em}
\setlist{itemsep=3pt,topsep=5pt,leftmargin=1.6em}
\newtheorem{theorem}{Theorem}[section]
\newtheorem{lemma}[theorem]{Lemma}
\newtheorem{proposition}[theorem]{Proposition}
\newtheorem{corollary}[theorem]{Corollary}
\theoremstyle{definition}\newtheorem{definition}[theorem]{Definition}
\theoremstyle{remark}\newtheorem{remark}[theorem]{Remark}
\newcommand{\ipra}{\mathsf{iPRA}}
\newcommand{\N}{\mathbb N}
\newcommand{\FV}{\operatorname{FV}}
\renewcommand{\paragraph}[1]{\par\medskip\noindent\textbf{#1}\enspace}
\newcommand{\file}[1]{\expandafter\nolinkurl\expandafter{\detokenize{#1}}}
\lstset{basicstyle=\ttfamily\footnotesize,columns=fullflexible,keepspaces=true,
 breaklines=true,showstringspaces=false,frame=single,rulecolor=\color{black!15},
 backgroundcolor=\color{pale},aboveskip=9pt,belowskip=9pt}
\pagestyle{fancy}\fancyhf{}
\fancyhead[L]{\small Equality from monus; neutral major premises}
\fancyhead[R]{\small Intuitionistic primitive recursive arithmetic}
\fancyfoot[C]{\thepage}
\hypersetup{pdftitle={Equality from Monus and Neutral Major Premises: Two Reductions for iPRA},
 pdfauthor={Research prepared for Vladimir Reshetnikov}}
\title{\color{ink}\bfseries Equality from Monus\\and Neutral Major Premises\\[8pt]
\large Two reductions for intuitionistic primitive recursive arithmetic}
\author{Research prepared for Vladimir Reshetnikov}
\date{10 November 2026}
\begin{document}
\maketitle
\begin{abstract}
We obtain two different reductions. First, the tied-seed arithmetic of the
bounded-initializer draft does not need a primitive Boolean test of actual
equality. Two equations for truncated subtraction suffice. A gated
complement identity proves equality reflection using only literal-equation
induction, with its passive parameter fixed. A bounded selector then
repairs the native logical bootstrap without assuming an unrestricted
numerical selector. Together with the explicitly attributed bounded
compiler, this gives a whole-domain bi-interpretation with full first-order
intuitionistic primitive recursive arithmetic.

Second, a separate proof protocol for the existing initialized-iterator
calculus uses a derived truth as the major premise of commands that
previously had none. It removes the checker's optional-major branch.
The induction mode is stored in the initializer field, whose falsity in
the introduction branch preserves the restriction to literal equations.
Complete implementations, precise input invariants, proof certificates,
source measurements, and an audited comparison accompany this result.

The mathematical basis and the measured checker are different
presentations. Their respective costs and novelty are not conflated.
The measured family improves the inspected byte, token, and syntax-tree
envelope, but different programs attain its coordinates. An absolute
smallest axiomatization and a single implementation below every old
coordinate have not been established. The written metatheorems have not
been verified in Lean or Rocq.
\end{abstract}
\begin{center}
\fcolorbox{ink}{pale}{\begin{minipage}{.92\linewidth}
\textbf{Exact status.} There is a new equality-reflection proof and a
smaller complete implementation of the inherited native calculus.
There is no basis for claiming the globally simplest theory or that the
new monus basis itself has the smallest checker. The requested broad
foundational and implementation breakthrough is therefore only partly
achieved.
\end{minipage}}
\end{center}
\clearpage
\setcounter{tocdepth}{1}
\begingroup\small\setlength{\parskip}{0pt}\tableofcontents\endgroup
\clearpage
\section{Results, comparison, and limits}
\label{sec:scope}

The supplied repository was inspected at commit
\begin{center}\small\file{76a2651601d17b34d46753f189ea026a21ce8d29}.\end{center}
Its \file{docs/drafts} directory contains 17 ZIP archives. The recursive
inventory, including embedded historical archives, contains 32 distinct
archive contents and 443 distinct Python source contents. Direct archive
bytes were checked against their Git blob identities. The snapshot,
recursive census, exact source identities, and current-checker replay
results are included in \file{evidence/}.\cite{snapshot}

The incumbent complete-checker envelope is \textbf{558 normalized bytes,
262 significant tokens, and 296 abstract-syntax-tree nodes}. Different
programs attain these values. All three leaders occur in the
certified-zero package, including its ancillary profiles; older report
headlines are not the current benchmark.\cite{certified}

\begin{center}\small
\begin{tabular}{lrrrr}
\toprule
Incumbent complete program & Bytes & Tokens & AST & Gzip-9\\
\midrule
Certified zero, explicit list test &566&264&308&351\\
Certified zero, compact list test &558&264&310&345\\
Certified zero, generic traversal &598&262&296&350\\
\bottomrule
\end{tabular}
\end{center}

The generic traversal has a documented qualification: nonempty malformed
strings rely on the finite CPython recursion limit for rejection. It is
not a structurally terminating parser on every possible JSON value in an
ideal unlimited-stack execution. The list-based profiles avoid that
qualification. All source and compression measurements include the
complete Python module, its import, JSON input, and acceptance driver.

\subsection{The mathematical reduction}
The bounded-initializer draft gives a tied-seed recursor $B$ and a
primitive Boolean equality test $\epsilon$. It explicitly asks whether
the test's range and reflection axioms can be replaced by short recursive
equations.\cite{bounded} We give such a replacement, using monus $d$:
\[
 d(x,0)=x,\qquad d(x,Sy)=P(d(x,y)).
\]
Here predecessor $P$ is a fixed $B$-term, not a new primitive. A literal
equation invariant proves that
\[
 \epsilon(x,y):=G\bigl(Z(d(x,y)),Z(d(y,x))\bigr)
\]
is a Boolean characteristic of actual equality throughout the domain.
The gate $G$ and zero test $Z$ are also $B$-terms. There is no appeal to
an extensional program oracle, a standard-number calculation, or
quantified induction in this bootstrap.

The corresponding native logical bootstrap is another necessary part of
the result. The usual unrestricted numerical selector cannot simply be
assumed in a tied-seed term language. We construct a selector whose first
output is Boolean. It supplies contextual case splits for equation
conclusions before equality reflection, and formula selection after
reflection. This closes the logical dependency cycle.

The bounded compiler itself is inherited. Its cap coherence and finite
quantifier-free graph method are reproduced with attribution in the
appendices. Together with recurrence uniqueness for the new $d$, they
give both whole-domain return interpretations. We do not label the
compiler or its almost-unary term analysis as new.

\subsection{The implementation reduction}
Separately, the initialized-iterator native calculus can replace
``optionally prove a major premise'' by ``always prove a major premise.''
When no major is needed, the premise is the empty-array representation of
the already derivable truth $0=0$. The context initially contains that
truth. A command still has four fields, but the case/progress mode moves
to a field that is necessarily false-valued when introduction is selected.

This changes the proof protocol, not the established numerical $I$
basis. The accompanying invariant rules out a false-valued numeric major
that would otherwise collide with the checker's closure-match sentinel.
The construction is therefore more than deleting an \texttt{if} statement:
it requires a conservative truth assumption, mode relocation, a
certificate translation, and a proof about every successful raw run.

\begin{center}\small
\begin{tabular}{lrrrr}
\toprule
Complete program & Bytes & Tokens & AST & Gzip-9\\
\midrule
Incumbent separate minima &558&262&296&---\\
New explicit list test &550&263&307&349\\
New byte profile &542&263&309&343\\
New generic traversal &578&261&295&346\\
\midrule
New separate minima &542&261&295&---\\
\bottomrule
\end{tabular}
\end{center}


The three new envelope coordinates are attained by different programs.
The explicit-list program is the clearest complete source; the byte
profile inherits an exact-JSON list-test idiom; the generic profile has
the same runtime qualification as its incumbent. No single program in
this report is below all three incumbent minima. The new monus theory
has not been assigned any of these checker sizes.

\subsection{What ``below the lower bound'' can mean here}
The draft sizes are achieved upper bounds on an unknown optimal
implementation size. They are not proved lower bounds on all possible
axiomatizations. We report exact improvements over their observed
coordinate envelope under the unchanged normalizer. A theorem of
absolute minimality would need a fixed description language and a
separate lower-bound proof.

The two results are useful, but they do not justify announcing the
globally simplest iPRA foundation or a large simultaneous implementation
breakthrough. The remaining gap is concrete: implement the new arithmetic
basis with all its logical and validation obligations, or find a single
plain checker that improves every current coordinate.

\subsection{Scope of iPRA and provenance}
Here iPRA is the \emph{full first-order intuitionistic} theory in a
language indexed by fixed primitive-recursive descriptions, with their
defining equations, successor boundary laws, and quantifier-free
induction. Atomic induction is an equivalent familiar formulation in the
full primitive-recursive language; Salehi uses that formulation.\cite{salehi}
We are not substituting the quantifier-free equational fragment for the
full intuitionistic theory.

The corpus audit found no earlier neutral-major protocol with the mode
relocation and live-context argument used here. It found an uncompleted
suggestion to add monus to the tied-seed basis, and an explicit question
about replacing equality reflection by recursion. Thus the new arithmetic
claim is the completed gated-complement and native-bootstrap proof, not
the first suggestion of the symbol $d$. These novelty statements concern
the inspected repository snapshot, not the entire mathematical literature.

All mathematical metatheorems below are written proofs. Executable
checks cover finite examples, certificate rejection, and source
accounting. They do not replace a Lean or Rocq proof of the metatheorems.

\section{Every primitive of the two native presentations}
\label{sec:rules}

The logical rules are shared. The numerical computation schemas differ,
and that difference is kept explicit.

\subsection{Numerical syntax}
The new arithmetic basis has finite terms
\[
 t::=x\mid0\mid St\mid d(t,t)\mid B(n;u,v.h).
\]
Only $u,v$ in $h$ are bound. Its remaining variables are passive parameters,
fixed during a recursion. Put $1=S0$. The two recursion schemas are
\begin{align}
 B(0;u,v.h)&=S h(0,0),\tag{B0}\label{basis:b0}\\
 B(Sn;u,v.h)&=h(n,B(n;u,v.h)).\tag{BS}\label{basis:bs}
\end{align}
The monus schemas are
\[
 d(x,0)=x,\qquad d(x,Sy)=P(d(x,y)),
 \tag{M0, MS}
\]
where
\[
 C(n,a)=B(n;u,v.a),\quad W(n)=B(n;u,v.u),\quad
 P(n)=W(C(n,n)).
\]
All substitutions are simultaneous where appropriate and avoid capture.

The independently measured checker keeps the inherited language
\[
 t::=x\mid0\mid St\mid I(n,a;v.h)
\]
and its equations
\[
 I(0,a;v.h)=a,\qquad I(Sn,a;v.h)=h(I(n,a;v.h)).
 \tag{I0, IS}
\]
The update receives its previous state, not the recursion counter. Its
direct computation packets have the certified-zero syntactic restrictions
specified in Section~\ref{sec:checker}; equality transport restores all
unrestricted instances. These $I$ equations are not axioms of the new
$B+d$ basis.

In both presentations, the binder notation abbreviates a staged ordinary
first-order signature. Each fixed finite earlier-stage body introduces
one function symbol whose arguments are its counter and passive
parameters, together with the independent initializer in the $I$ case.
Bodies are never quantified over. Every term and proof uses finitely many
stages. There is no internal uniformly total evaluator of arbitrary
primitive-recursive codes.

\subsection{Formula constructors and readings}
Native formulas are generated by numerical equations and the paired
binder:
\[
 A::=(s=t)\mid D_{x,y}A.
\]
The auxiliary closure $Q_yA$ may be used in certificates and contexts,
or admitted as a definable formula constructor. Its representative is
$D_{a,y}A$ with unused $a$. Their ordinary readings are
\[
 Q_yA(y):=A(0)\to\forall y\,A(y),\qquad
 D_{x,y}A(x,y):=\exists x\bigl(A(x,0)\to\forall y\,A(x,y)\bigr).
\]
These readings explain the interpretation; ordinary logical rules are
not silently included among the native rules. The designated native
contradiction is $0=1$.

Contexts $\Gamma$ are finite reusable lists. All logical rules are:
\begin{align*}
 &\frac{}{\Gamma\vdash A}\ (A\in\Gamma) &&\text{(Hyp)},\\[4pt]
 &\frac{\Gamma\vdash s=t\quad\Gamma\vdash F(s)}
        {\Gamma\vdash F(t)} &&\text{(Tr)},\\[4pt]
 &\frac{\Gamma\vdash Q_yA(y)\quad\Gamma\vdash A(0)}
        {\Gamma\vdash A(t)} &&\text{(QE)},\\[4pt]
 &\frac{\Gamma\vdash D_{x,y}A(x,y)\quad
         \Gamma,Q_yA(b,y)\vdash F}
        {\Gamma\vdash F} &&\text{(DE)},\\[4pt]
 &\frac{\Gamma,A(t,0)\vdash A(t,Sz)}
        {\Gamma\vdash D_{x,y}A(x,y)} &&\text{(DI-case)},\\[4pt]
 &\frac{\Gamma,E(t,z)\vdash E(t,Sz)}
        {\Gamma\vdash D_{x,y}E(x,y)} &&\text{(DI-progress)}.
\end{align*}
In Tr, $F$ is any fixed native formula template. In DE, $b$ is fresh for
the retained context, conclusion, and fixed matrix parameters. In both
introduction rules, $z$ is fresh for the context, witness, and passive
matrix parameters; the displayed assumption is discharged.
\textbf{DI-progress is restricted to a literal numerical equation.}
An equation beneath a binder does not satisfy that restriction.

The numerical computation schemas are axiom leaves. Equivalently, the
implementation uses
\[
 \frac{\Gamma,H\vdash A}{\Gamma\vdash A}
 \qquad(H\text{ a permitted computation equation}).\tag{Supply}
\]
There are no additional native logical axioms, successor axioms, equality
axioms, or induction schemas. In particular, reflexivity and explosion
are derived below, not omitted costs hidden in a conventional logic.

Weakening, exchange, contraction, renaming, numerical substitution, and
cut of an ordinary assumption are admissible. Their proofs are induction
on finite derivations; eigenparameters are renamed before substitution
or cut insertion. A numerical substitution leaves a literal equation
literal. For syntactically restricted $I$ computations, the certified-zero
repair is applied whenever substitution creates a forbidden instance.

\subsection{The ordinary comparison theory, with its entire logic}
The next section first proves the arithmetic result in a conventional
presentation $T_{Bd}$, to make the invariant easy to inspect. Its
formula grammar is
\[
 A::=(s=t)\mid\bot\mid A\land A\mid A\lor A\mid A\to A
          \mid\forall xA\mid\exists xA.
\]
Its logical rules are exactly intuitionistic natural deduction:

\begin{center}\small
\begin{tabularx}{\linewidth}{@{}lX@{}}
\toprule
Form & Rules\\\midrule
Assumptions & Use any member of the reusable context; weakening, exchange,
and contraction.\\[3pt]
$\bot$ & From $\bot$ infer any formula; no introduction rule.\\[3pt]
$A\land B$ & Introduce from both conjuncts; eliminate to either conjunct.\\[3pt]
$A\lor B$ & Introduce from either disjunct; eliminate by two derivations of
the same conclusion, under the respective disjuncts, discharging them.\\[3pt]
$A\to B$ & Introduce by discharging assumption $A$ from a derivation of $B$;
eliminate with a proof of $A$.\\[3pt]
$\forall xA$ & Introduce from $A(z)$ with fresh $z$ absent from the retained
context and fixed parameters; eliminate by capture-free term substitution.\\[3pt]
$\exists xA$ & Introduce from $A(t)$; eliminate under $A(z)$ with $z$ fresh
for the context, conclusion, and fixed parameters, discharging $A(z)$.\\[3pt]
$=$ & Reflexivity $t=t$; from $s=t$ and $F(s)$ infer $F(t)$ for every
fixed formula template, with capture avoidance.\\
\bottomrule
\end{tabularx}
\end{center}

Its additional arithmetic principles are exactly the four displayed
computation schemas, coverage, separation, and literal-equation induction
listed next. The later native bootstrap derives the conventional logical
interface rather than presupposing it.

Successor is eliminable from the native numerical signature:
$S(a)=C(0,a)$. One can use this abbreviation everywhere, including the
four computation schemas. Keeping $S$ visible makes the proofs easier
to read. This elimination is a definitional equivalence, not an additional
measured checker reduction.

% Includable proof fragment. Requires amsmath, amssymb, amsthm, mathtools.
% Environments theorem, lemma, proposition, corollary assumed.
% This is a new derivation relative to the inspected current draft corpus.
\section{Replacing primitive numerical equality by two monus equations}
\label{sec:monus-basis}

We work initially in ordinary intuitionistic first-order logic with actual
numerical equality. The exact numerical language consists of $0,S$, the
staged tied-seed constructor $B$, and a binary symbol $d$. For each fixed
finite earlier-stage body $h(u,v,\vec p)$ the two $B$ schemas are
\[
 B_h(0,\vec p)=S h(0,0,\vec p),\qquad
 B_h(Sn,\vec p)=h(n,B_h(n,\vec p),\vec p).
 \tag{B}
\]
The two additional computation schemas are
\[
 d(x,0)=x,\qquad d(x,Sy)=P(d(x,y)).
 \tag{M}\label{monus:M}
\]
Here $P$ is the fixed $B$-term constructed below; it is not a further
primitive. Numerical coverage and separation are explicit:
\[
 \forall x\,(x=0\lor\exists y\,x=Sy),\qquad 0=S0\to\bot.
 \tag{CS}\label{monus:CS}
\]
The sole induction schema is contextual induction on a literal equation:
\[
 \frac{\Gamma\vdash E(0,\vec p)\quad
       \Gamma,E(z,\vec p)\vdash E(Sz,\vec p)}
      {\Gamma\vdash E(t,\vec p)},
 \tag{EI}\label{monus:EI}
\]
with $z$ fresh for $\Gamma,t,\vec p$. All schemas are universally closed
in their free numerical variables. There are no function variables or
uniform evaluation symbols. Call this ordinary first-order theory $T_{Bd}$.

Coverage is an essential initial resource in the proof below. It must not
be claimed as a consequence of the positive computation equations and
literal-equation induction: the coordinatewise square of the standard
model satisfies those positive equations and their equation-induction
instances, while numerical coverage fails. A native calculus with a sound
numerical case rule may supply this resource instead, but its reconstruction
of ordinary logical rules must be justified independently.

\subsection{Elementary terms available before equality reflection}
Put $1=S0$ and define
\[
 C(n,a)=B(n;u,v.a),\qquad W(n)=B(n;u,v.u),\qquad
 P(n)=W(C(n,n)),\qquad Z(n)=C(n,0).
\]
The computation schemas alone give
\[
 C(0,a)=Sa,\quad C(Sn,a)=a,\quad W(0)=1,\quad W(Sn)=n,
\]
and therefore
\[
 P(0)=0,\quad P(Sn)=n,\qquad Z(0)=1,\quad Z(Sn)=0.
 \tag{PZ}\label{monus:PZ}
\]
Successor is injective, by application of $P$. It never equals zero:
$Sx=0$ would imply $x=P(0)=0$, and hence $S0=0$, contradicting
separation. Coverage and the endpoints of $Z$ now prove
\[
 Z(x)=0\lor Z(x)=1,\qquad Z(x)=1\ \longleftrightarrow\ x=0.
 \tag{Z}\label{monus:Z}
\]
These are proofs in intuitionistic logic using the explicit coverage axiom;
they do not use induction on a disjunction.

Define the one-sided gate
\[
 G(n,a)=P\bigl(P(B(a;u,v.C(Z(n),v)))\bigr).
\]
Literal-equation induction on $a$ proves
\[
 G(0,a)=0,\qquad G(Sn,a)=a.
 \tag{G}\label{monus:G}
\]
For the zero control the inner iterator has initial value $1$ and identity
update, so it is always $1$. For a successor control its initial value is
$2$ and its update is successor, so its value at $a$ is $S(Sa)$. Applying
$P$ twice gives the two equations. In particular, $G(1,a)=a$.

The inherited zero-start recursion macro is also available without
numerical equality:
\[
 L_h(n,\vec p)=B\bigl(Sn;r,s.G(r,h(Pr,s,\vec p))\bigr).
\]
Its first actual update is zero, and the subsequent updates are $h$;
therefore its equations are
\[
 L_h(0,\vec p)=0,\qquad
 L_h(Sn,\vec p)=h(n,L_h(n,\vec p),\vec p).
 \tag{L}
\]
No arbitrary initialization has been assumed here.

\subsection{Four preliminary monus identities}

\begin{lemma}\label{monus:preliminary}
The theory proves, using only literal-equation induction,
\begin{align}
 d(0,y)&=0,\label{monus:leftzero}\\
 P(d(Sx,y))&=d(x,y),\label{monus:commute}\\
 d(Sx,Sy)&=d(x,y),\label{monus:diagonal}\\
 d(x,x)&=0,\qquad d(Sx,x)=1.\label{monus:diagonal-values}
\end{align}
\end{lemma}
\begin{proof}
For \eqref{monus:leftzero}, induct on $y$ using $P0=0$.
For \eqref{monus:commute}, keep $x$ fixed and induct on $y$.
The base is $P(Sx)=x$. At the step,
\[
 P(d(Sx,Sy))=P(P(d(Sx,y)))
           =P(d(x,y))=d(x,Sy),
\]
where the middle equality applies $P$ to the induction equation.
Equation \eqref{monus:diagonal} follows immediately from
\eqref{monus:M} and \eqref{monus:commute}.
Finally induct on $x$ in the two literal equations of
\eqref{monus:diagonal-values}. Their bases are $d(0,0)=0$ and
$d(1,0)=1$, and their steps are exactly \eqref{monus:diagonal}.
No induction hypothesis is generalized or specialized to a changed
passive parameter.
\end{proof}

\subsection{The gated complement identity}
The key step is an equation invariant that permits the ordinary
conditional complement argument before quantifier-free induction is
available.

\begin{lemma}[Gated complement]\label{monus:complement}
For all $x,y$ the theory proves
\[
 G\bigl(Z(d(y,x)),d(x,d(x,y))\bigr)
       =G\bigl(Z(d(y,x)),y\bigr).
 \tag{GC}\label{monus:GC}
\]
Consequently,
\[
 d(y,x)=0\quad\longrightarrow\quad d(x,d(x,y))=y.
 \tag{C}\label{monus:C}
\]
\end{lemma}
\begin{proof}
Keep $x$ fixed and apply literal-equation induction to
\eqref{monus:GC} as an equation in $y$. At $y=0$, the gate control
is $Z(d(0,x))=1$, and both output arguments are zero by
$d(x,0)=x$ and $d(x,x)=0$.

For the step, retain the induction equation and set
$a=d(Sy,x)$. Use coverage on $a$. If $a$ is a successor, $Z(a)=0$
and both sides of the new equation are zero by \eqref{monus:G}.
It remains to consider $a=0$.

By \eqref{monus:commute},
\[
 d(y,x)=P(d(Sy,x))=0.
\]
Thus the induction equation, with both gates simplified at control $1$,
gives
\[
 d(x,t)=y,\qquad t=d(x,y).
 \tag{IH}\label{monus:IH}
\]
If $t=0$, then $d(x,t)=x$, so $x=y$. But this makes
$a=d(Sy,y)=1$, contradicting $a=0$. Hence $t\ne0$.
Coverage therefore provides $u$ with $t=Su$. The monus successor
schema identifies
\[
 u=P(t)=d(x,Sy).
\]
Put $v=d(x,u)$. Equation \eqref{monus:IH} and the successor schema give
\[
 P(v)=P(d(x,u))=d(x,Su)=d(x,t)=y.
\]
If $v=0$, then $y=0$. Hence $t=d(x,0)=x=Su$, and
\[
 v=d(x,u)=d(Su,u)=1,
\]
a contradiction. Thus $v\ne0$, and coverage gives $v=Sz$.
Since $P(v)=y$, the predecessor equations give $z=y$, so $v=Sy$.
Together with $u=d(x,Sy)$ this is
\[
 d(x,d(x,Sy))=Sy,
\]
which is exactly the required step with gate control $Z(a)=1$.
All contradictions use only $0\ne1$ and intuitionistic falsity elimination.
All case splits occur inside the step derivation with the original
induction equation retained. The only induction matrix used was the
single literal equation \eqref{monus:GC}.

Finally if $d(y,x)=0$, the control of both gates in
\eqref{monus:GC} is $1$, giving \eqref{monus:C}.
\end{proof}

\begin{corollary}[Antisymmetry before order theory]\label{monus:antisymmetry}
The theory proves
\[
 d(x,y)=0\ \wedge\ d(y,x)=0\quad\longrightarrow\quad x=y.
\]
\end{corollary}
\begin{proof}
Equation \eqref{monus:C} gives $d(x,d(x,y))=y$. Substituting
$d(x,y)=0$ and then $d(x,0)=x$ gives $x=y$.
\end{proof}

\subsection{A characteristic of actual equality}
Define the fixed numerical term
\[
 \epsilon(x,y)=G\bigl(Z(d(x,y)),Z(d(y,x))\bigr).
 \tag{E}\label{monus:E}
\]

\begin{theorem}[Whole-domain equality reflection]\label{monus:reflection}
The theory $T_{Bd}$ proves
\[
 \epsilon(x,y)=0\lor\epsilon(x,y)=1,
 \qquad
 \epsilon(x,y)=1\ \longleftrightarrow\ x=y.
\]
These statements concern actual equality throughout the domain.
\end{theorem}
\begin{proof}
Each input of the gate is zero or one by \eqref{monus:Z}. The two
gate equations prove the displayed range. They also prove
\[
 \epsilon(x,y)=1\quad\longleftrightarrow\quad
                  d(x,y)=0\wedge d(y,x)=0:
\]
if the first gate input were zero, the output would be zero; when it
is one, the output is the second input. Use $0\ne1$ and
\eqref{monus:Z} in both directions. Antisymmetry proves the right-to-left
numerical identification required for reflection. Conversely $x=y$
implies both differences are zero by $d(x,x)=0$ and equality substitution,
so the gate output is one.
\end{proof}

In particular numerical equality is decidable, but excluded middle for
arbitrary quantified formulas has not been postulated. Every fixed
quantifier-free formula has a Boolean characteristic term: use
\eqref{monus:E} for atoms, zero for falsity, $Z$ for negation, and $G$
for conjunction; disjunction and implication follow from the Boolean
identities. Implication is legitimate intuitionistically because the
antecedent is decidable. Literal-equation induction on the characteristic
being one yields contextual quantifier-free induction.

\subsection{Consequences for the existing bounded compiler}

The equality term just proved satisfies precisely the range and actual
reflection axioms of the inspected $B+\epsilon$ bounded-initializer
presentation. Its finite descriptions and literal-equation induction
instances translate by substitution of \eqref{monus:E}. Therefore the
existing bounded-initialization construction, cap-coherence argument,
and finite quantifier-free function graphs apply to $T_{Bd}$.
They give a whole-domain interpretation of iPRA with actual equality.

In the reverse direction, iPRA interprets each fixed $B_h$ by primitive
recursion with base $S h(0,0,\vec p)$ and step $h$, and interprets $d$
by ordinary truncated subtraction. It proves the two monus equations,
coverage, separation, and every source equation-induction instance.
The common numerical domain is unchanged.

For the source-side composite, first identify the $d$-free constant-body
$C$ and counter-body $W$ recursors by recurrence uniqueness, and hence
their predecessor macro $P$. The imported primitive-recursive monus
graph is provably equal to the original $d$: induct on its second
argument using its defining equations and \eqref{monus:M}.
When a graph-defined function is used in this comparison, expand its
finite quantifier-free graph cases; the quantifier-free induction already
derived above licenses this definitional-extension argument. For each
fixed staged $B$ symbol, recurrence uniqueness similarly identifies the
composite with the original symbol, proceeding by external induction on
the finite description. The iPRA-side composite follows from primitive-
recursion uniqueness in the same way. Thus the inherited compiler proof
extends to full-domain bi-interpretability and equiconsistency.

This is a new arithmetic presentation with a transparent reflection
bootstrap. It is not, by itself, a proof-checker size reduction. Its
ordinary logical rules, coverage, separation, all four computation
schemas, and equation induction must be included in any honest
implementation comparison. A checker for a different inherited calculus
cannot be assigned to this new theory merely because both are
bi-interpretable with iPRA.

% Native Eq/Q/D bootstrap for the new B+monus arithmetic basis.
% Requires the core native rules to have already been explicitly stated.
\section{Native cases before a numerical equality characteristic}
\label{sec:monus-native-bootstrap}

There is a potential circularity in importing the preceding arithmetic
proof directly into the native Eq/$Q$/$D$ calculus. Its theorem-level
numerical case rule is not automatically a contextual split on $t=0$
and $t=Sz$ when the retained induction hypothesis mentions $t$. Moreover,
the tied-seed arithmetic has no unrestricted numerical selector among its
term functions. The following bounded selector removes that circularity.
Only hypothesis use, full formula transport, the native $Q/D$ rules,
the $B$ computation schemas, and literal-equation progress are used.
No ordinary implication, conjunction, disjunction, or quantifier rule is
assumed in this section.

\subsection{Initial native resources}
The native rules derive numerical cases and literal-equation induction
with their usual freshness conditions, using an unused witness in $D$,
then $D$ elimination and $Q$ specialization. The computation
$C(S0,p)=p$, used as both premises of transport with motive $w=p$,
gives $p=p$. Symmetry, transitivity, and numerical congruence follow
from the usual transport motives. The $B$ macros $P,Z,G$ of the previous
section are therefore available by their same equation derivations.

Equation explosion is also available. Under $0=1$, transport between
$C(0,a)=Sa$ and $C(1,a)=a$ gives $Sa=a$ for arbitrary $a$.
Literal-equation induction on $a=0$ then proves every $a=0$ in this
retained context. Hence every numerical equation follows. External
induction on a fixed native formula, using case introduction for each
$D$ and the unused-witness representation of $Q$, gives explosion for
all native formulas. This is an external finite proof transformation;
it is not compound-formula induction in the object theory.

For a numerical term $b$, say $\mathrm{FR}_{\{0,1\}}(b)$ holds when,
for every fixed native formula template $K$, the following metarule is
available in an arbitrary retained context:
\[
 \Gamma\vdash K(0),\quad\Gamma\vdash K(1)
       \quad\Longrightarrow\quad\Gamma\vdash K(b).
 \tag{FR}\label{monus:native-FR}
\]
The context may mention $b$; the premises add no equations $b=0$ or
$b=1$. Native numerical cases on the fresh template $w\mapsto K(Zw)$,
using $Z0=1$ and $Z(Sw)=0$, prove
$\mathrm{FR}_{\{0,1\}}(Zt)$ for every $t$. In particular,
$Z(Zb)=b$ follows for every such finite-range $b$ by applying FR to
that equation's two numeral instances.

\subsection{A selector whose first output is Boolean}
Define
\[
 J(a,b)=C(Zb,a),\qquad
 H(y,b,u)=J\bigl(G(y,u),G(Zy,b)\bigr).
 \tag{H}\label{monus:H}
\]
Direct computations give
\[
 J(a,0)=a,\quad J(a,1)=Sa,\qquad
 H(0,b,u)=b\quad(b\in\{0,1\}),\quad H(Sy,b,u)=u.
 \tag{H endpoints}\label{monus:H-endpoints}
\]
The first endpoint also holds for any term $b$ with the displayed
functional range rule: apply that rule to the fixed equation template.
The successor endpoint holds without any restriction on $b$.

This is not an unrestricted numerical selector. Its zero output is
Boolean; its successor output is arbitrary. Consequently its existence
does not contradict the one-coordinate majorant bound for tied-seed
terms. It is exactly enough for contextual Boolean case elimination.

\begin{lemma}[Contextual Boolean cases for equation conclusions]
\label{monus:contextual-Boolean}
Suppose $\mathrm{FR}_{\{0,1\}}(b)$ holds. If
\[
 \Gamma,b=0\vdash U=V,\qquad
 \Gamma,b=1\vdash U=V,
\]
then $\Gamma\vdash U=V$ is derivable in the native calculus. The retained
context and the conclusion may mention $b$.
\end{lemma}
\begin{proof}
With fresh $a,y$, define
\[
 K(x)=D_{a,y}\bigl(H(y,x,U)=H(y,b,V)\bigr),
\]
where $a$ is unused. For $r=0,1$, use case introduction to prove $K(r)$.
Its local zero assumption reduces to $r=b$, by the Boolean endpoint
of $H$ for $r$ and for $b$. Symmetry gives $b=r$, so cut in the
corresponding supplied branch proof. The successor matrix is $U=V$
by the other endpoint of $H$. This constructs $\Gamma\vdash K(r)$
for both numeral values, discharging the local equation.

Apply FR to obtain $\Gamma\vdash K(b)$. Paired elimination opens its
conditional closure. Its zero instance is $b=b$, supplied by
reflexivity and the endpoint computations. Specialization at $1$
returns $U=V$. All temporary variables are fresh, and all temporary
assumptions have been discharged.
\end{proof}

\subsection{Contextual zero/successor splits for arbitrary terms}
Native numerical cases, applied to fresh variables in the following
literal-equation templates, prove
\[
 G(Zt,t)=0,
 \qquad
 G(Z(Zt),t)=G(Z(Zt),S(Pt)).
 \tag{Reconstruction}\label{monus:native-reconstruction}
\]
For the first identity the zero branch is $G(1,0)=0$ and every
successor branch has control zero. For the second, the zero branch
has control zero, while the successor branch has control one and
$P(Sx)=x$. There is no substitution into a retained context.

\begin{lemma}[Contextual numerical cases for equations]
\label{monus:contextual-numerical}
Let $z$ be fresh for $\Gamma,t,U,V$. If
\[
 \Gamma,t=0\vdash U=V,\qquad
 \Gamma,t=Sz\vdash U=V,
\]
then $\Gamma\vdash U=V$ is derivable natively.
\end{lemma}
\begin{proof}
Apply Lemma~\ref{monus:contextual-Boolean} to $b=Zt$.
Under $Zt=1$, the first reconstruction identity yields $t=0$,
so the first given branch applies. Under $Zt=0$, the second identity
has control $Z0=1$, so it yields $t=S(Pt)$. Substitute $Pt$ for the
fresh $z$ in the second given branch and cut in this equation.
The conclusion is the same equation $U=V$ in both Boolean branches.
\end{proof}

Every case split in the gated-complement proof has an equation as its
ultimate conclusion. The zero cases that contradict $0\ne1$ are
replaced in the native calculus by the equation-explosion transformation
above. Thus the proof of the gated-complement invariant, antisymmetry,
and the two directions of equality reflection translates natively
without importing ordinary logical rules. Its only induction matrix
remains the single equation in the displayed invariant.

\subsection{Reconstructing ordinary logic after reflection}
The equality term $\epsilon(s,t)$ now has the native functional Boolean
range rule and satisfies
\[
 s=t\ \dashv\vdash\ \epsilon(s,t)=1
\]
in every retained context. Its finite range follows from the two
functional range rules for its $Z$ inputs and the gate endpoints.
For an atomic formula $E\equiv(s=t)$ define a gate by
\[
 \mathcal G_v(E)\;:\equiv\;
        G(v,\epsilon(s,t))=G(v,1).
\]
It is provable at $v=0$ and interderivable with $E$ at each $v=Sw$.
For atomic $E\equiv(s=t)$ and $F\equiv(r=w)$ define a formula selector
\[
 \mathcal M_v(E,F)\;:\equiv\;
        H(v,\epsilon(s,t),\epsilon(r,w))=1.
\]
The Boolean endpoint of $H$ gives $\mathcal M_0(E,F)\dashv\vdash E$;
its successor endpoint gives $\mathcal M_{Sz}(E,F)\dashv\vdash F$.
The zero branch is bounded, so no unrestricted numerical selector has
been introduced.

Extend the gate through $D$, replace $Q$ by its unused-witness $D$
representation, and align binder depths using the usual gate padding.
Then extend the formula selector through matched $D$ binders by external
recursion on their finite syntax. Binder congruence uses only the native
case rule, paired elimination, transport, and ordinary-assumption cut,
so it is already available. These are the same structural logical
constructions as in the inherited Choice-DI bridge~\cite{certified}, with
the atomic selector just given replacing its unbounded numerical selector.
Both inherited derivations are included under
\file{reference/initialized_original/} as
\file{inherited_choice_basis.tex} and
\file{inherited_ordinary_bridge.tex}.

For explicitness, put $p(z)=L(z;u,v.Zv)$, so $p(1)=1$ and
$p(2)=0$. With all unused coordinates fresh, the ordinary macros are
\begin{align*}
 A\Rightarrow_*B&=D_{a,z}\mathcal M_z(A,B),\\
 A\wedge_*B&=D_{a,z}\mathcal G_z\mathcal M_{p(z)}(A,B),\\
 \exists_*x\,A(x)&=D_{x,y}\mathcal G_y A(x),\\
 A\vee_*B&=\exists_*x\,\mathcal M_x(A,B),\\
 \forall_*x\,A(x)&=A(0)\wedge_*D_{a,z}\mathcal G_z A(z),\\
 \bot_*&=(0=1).
\end{align*}
Implication introduction uses the matrix's zero assumption; elimination
opens its closure and specializes at one. The gate's provable zero
instance supplies conjunction and existential introduction, while
specialization at one or two supplies their eliminations. After implication
is available, ordinary contextual numerical cases for arbitrary formulas
follow by applying Native NC to an implication template; these cases
justify disjunction elimination and universal elimination. All freshness
and capture-avoidance conditions are the inherited ones.

These standard native macros for implication, conjunction, disjunction,
existence, universal quantification, and falsity therefore have their usual
intuitionistic introduction and elimination rules. Both formula
composites follow from the same gate and selector endpoint proofs.
After this stage the ordinary $T_{Bd}$ argument and the inherited
bounded compiler may be used without a logical bootstrap cycle.

\section{Both interpretations and their composites}
\label{sec:interpretations}

The preceding argument supplies an actual numerical equality
characteristic and full intuitionistic logic. The rest of the arithmetic
translation uses the bounded-initializer construction, reproduced in
Appendix~\ref{app:bounded}. This section explains exactly what is inherited
and what must additionally be checked for the new primitive monus.

\subsection{Why graphs, rather than unrestricted output terms}
Tied-seed terms can still be bounded by a primitive-recursive unary
function of one input coordinate. That does not prevent a function from
being definable by a finite choice of terms. The bounded compiler supplies,
for every fixed primitive-recursive description $f$, source terms
\[
 \tau_f(b,\vec x),\qquad M_f(b),
\]
with proved bounds and coherence: whenever both $b$ and $c$ bound the
arguments, $\tau_f(b,\vec x)=\tau_f(c,\vec x)$. Neither $f$ nor a
description code is quantified over inside the source theory.

For a nonempty tuple $(x_1,\ldots,x_k)$, let $C_i(\vec x)$ say that $i$
is its first largest coordinate in the internally derived discrete
order. These are finitely many exhaustive, disjoint, quantifier-free
cases. The graph
\[
 \Phi_f(\vec x,z)=\bigvee_{i=1}^k
       \bigl(C_i(\vec x)\land z=\tau_f(x_i,\vec x)\bigr)
\]
is provably total and single-valued. Its outputs are explicit terms on
each case. Nullary functions have a single fixed value term.

Finite value cases are crucial for induction preservation. Expanding
every occurrence of a graph-defined function in a finite term gives a
finite family of quantifier-free guards and source value terms. Refining
these cases for every term in a quantifier-free formula gives an
equivalent source quantifier-free formula. This applies to implication
and negation as well, because the guards cover all inputs and the values
agree on overlaps. Source quantifier-free induction, already derived
from literal-equation induction and our equality characteristic,
therefore proves every translated iPRA induction instance.

An arbitrary existential graph would not establish this fact. The finite
term-value construction in the appendix is the relevant part of the
inherited proof.

\subsection{Arithmetic round trips}
Interpret the new $B+d$ language in iPRA by ordinary primitive recursion:
the initial value for $B_h$ is $S h(0,0,\vec p)$, its update is $h$, and
$d$ is ordinary truncated subtraction. The two $B$ schemas and the two
monus schemas hold. Interpret native equations as actual equality and
$Q,D$ by their ordinary readings. Coverage proves DI-case and equation
induction proves DI-progress. Every native rule is therefore sound.

Interpret iPRA back by the finite graphs above. The computation and
finite-case lemmas prove the defining equations and all induction
instances, and predecessor plus separation prove the successor boundary
laws. Both numerical interpretations have domain formula $x=x$ and use
actual equality.

On the iPRA side, external induction on each fixed primitive-recursive
description, with internal quantifier-free induction at a recursion
node, identifies its composite with the original function.

On the source side, first preserve the $d$-free constant-body recursor
$C$ and counter-body recursor $W$ by their equations and recurrence
uniqueness. This preserves $P=W(C(n,n))$. Next preserve primitive $d$
using its two equations, the preserved $P$, and finite quantifier-free
graph elimination. The constructed difference
$\widehat d$ of the inherited bridge has the same two recursive equations
as our primitive $d$. Literal-equation induction in its second argument
proves $\widehat d(x,y)=d(x,y)$; this source-level proof is included at the
end of the order appendix. Thus the new primitive is recovered, not just
some externally equivalent standard function. The equality characteristic
is then identified by its proved Boolean range and reflection. For each fixed staged
$B_h$, external induction on its earlier-stage body identifies the two
updates, and equation induction on the counter identifies their results.
When a graph-defined function occurs in a comparison, its finite
quantifier-free cases justify the induction in the original language.

\subsection{Formula round trips}
The gate and selector endpoint proofs identify each ordinary connective
with its native macro. The nontrivial native composite is
\[
 D_{x,y}A(x,y)\quad\dashv\vdash\quad
 \exists_*x\bigl(A(x,0)\Rightarrow_*\forall_*y\,A(x,y)\bigr).
\]
Forward, open the paired formula at a fresh witness. Under its zero
instance, QE provides every instance, giving the universal and then the
implication and existential macros. Reverse, eliminate the translated
existential, retain the matrix's zero instance, and use its implication
and universal at $Sz$. Case introduction reconstructs the original
paired formula. The unused-witness representative gives the analogous
composite for $Q$.

External structural induction on a fixed formula now proves both
formula composites, including arbitrary quantified formulas. All
quantifiers range over the original whole domain. The identity relation
$x=y$ defines each numerical isomorphism and preserves every operation,
relation, and translated formula. Ordinary falsity corresponds to the
native contradiction $0=1$, whose explosion was explicitly derived.

\begin{theorem}[The new native arithmetic basis]
The native Eq/$Q$/$D$ theory with B0, BS, M0, and MS, and exactly the
rules of Section~\ref{sec:rules}, is whole-domain bi-interpretable and
equiconsistent with full first-order iPRA. Both composite domain
isomorphisms are defined by actual equality.
\end{theorem}

The equiconsistency conclusion follows from effective proof
translations preserving contradiction. It is not a claim that either
theory proves its own consistency. Nor does it identify the implementation
cost of this basis with the smaller, independently measured $I$ checker.

\section{The smaller checker and its complete boundary}
\label{sec:checker}

This section concerns the initialized-iterator presentation specified in
Section~\ref{sec:rules}. It does not implement the new $B+d$ grammar.
The mathematical theorem is identity preservation of its native
consequence relation under a changed proof protocol. Thus its existing
whole-domain iPRA interpretations transfer unchanged.

\subsection{A precise sorted representation}
The inherited encoding is sorted: the same array can have different
meanings in a numerical position and a formula position. Its native
meanings are:
\begin{center}\small
\begin{tabular}{cll}
\toprule
Array arity & Numerical meaning & Formula meaning\\\midrule
0 & $0$ & $0=0$\\
1 & Successor of its child & Paired binder $D$\\
2 & $0$, annotations ignored & Equality of its two terms\\
3 & Initialized scalar iteration & Closure $Q$, last child is matrix\\
$\ge4$ & $0$, annotations ignored & $0=0$, annotations ignored\\
\bottomrule
\end{tabular}
\end{center}
Canonical paired formulas are \texttt{[B]}, canonical equations are
\texttt{[s,t]}, and canonical closures are \texttt{[B,B,B]}. The three
copies let list multiplication export exactly a closure when its major
is unary; all other nonempty array majors export a high-arity truth.
This arity construction is inherited, not a new contribution.\cite{arity}

The opener increases depth by $5-r$ at an array of arity $r$. Numerical
root depth is three; formula root depth is zero. Under $p$ protected
numerical slots, the corresponding depths are $3+2p$ and $2p$. Root
coordinates $0,1,2$ denote numerical zero. Coordinates $3+2j,4+2j$ denote
the $j$th protected slot. Negative admitted labels name rigid parameters.
A unary successor's child is decoded beneath two dummy zero slots;
the counter and initializer of an iterator are decoded beneath one dummy
zero slot, while its update is under its actual state slot. This is how
the common depth law implements the different sorted constructors.

At global floor $k$, a scalar in $[k,d)$ is retained. A scalar escaping
at numerical depth is the opening hole; its supplied image is checked at
numerical root depth three. An escaping scalar at formula depth erases
to the truth padding. A missing image at a genuinely used numerical
hole rejects. The proof that this implements capture-free opening is the
inherited structural induction over the sorted decoder. An image has
only rigid negative free labels, so copying it under protected slots
does not capture those free parameters. Large ignored arities retain
their constant meanings; they are not assumed to be idempotently
revalidated as raw trees.

Every non-hypothesis continuation decreases $k$. Incoming goals and
premises have been checked at the older floor, so a newly admitted
parameter is fresh for the retained context and conclusion. If a finite
raw certificate admits fractional labels in a newly opened interval,
they are distinct fresh parameters; all are included in a later
zero-instance substitution when needed.

\subsection{The restricted direct computation interface}
The initializer field must be nonempty to select computation. A false
predecessor selects the zero equation; a true predecessor selects its
successor equation. Thus the canonical directly supplied instances are
\[
 I(0,a;v.h)=a\quad(a\text{ not literal }0),
\]
\[
 I(Sn,a;v.h)=h(I(n,a;v.h))
 \quad(n,a\text{ not literal }0).
\]
This is a syntactic restriction, not a numerical nonzero test. The
inherited normalization is short and noncircular. Set
\[
 Z_*:=I(S(S0),S0;v.0).
\]
The permitted successor equation immediately gives $Z_*=0$: its
predecessor and seed are both the nonempty numeral one. For an excluded
computation, replace every offending predecessor or seed position
simultaneously by a fresh placeholder $w$ in a fixed equation motive
$E(w)$. Keep unrelated zeros in the update unchanged. Then $E(Z_*)$ is
a permitted computation, and transport along $Z_*=0$ proves $E(0)$.
The repair consists of one transport node and its two computation
premises. All unrestricted I0 and IS instances follow.\cite{certified}

Generated raw hypotheses need only have a known zero instance. For
fresh $\vec b$ absent from $\Gamma,A$, the rule
\[
 \vdash H(\vec0),\quad\Gamma,H(\vec b)\vdash A
 \quad\Longrightarrow\quad\Gamma\vdash A
\]
is admissible: substitute zero for the fresh names in the second
derivation, rename subsequent eigenparameters first, and cut in the
first proof. The modified computation mode still uses this rule.
It would be false to say every generated open hypothesis is uniformly
true in its newly admitted parameters.

\subsection{All fields and all command branches}
The input is one JSON array. Its last element is the requested goal.
Earlier four-field packets $[m,s,u,t]$ are consumed from the end.
The initial assumption is the empty array, a representation of the
derived theorem $0=0$. A hypothesis closes without consuming a packet.
Every other command first proves $m$ in the current context.

\begin{center}\small
\begin{tabularx}{\linewidth}{@{}lX@{}}
\toprule
Case & Meaning\\\midrule
Truthy $m$ & Check the substantive major. The old equality/closure endpoint
match can change the goal from $F(u)$ to $F(s)$. Export $m*3$ as the
paired elimination resource or a constant truth alias.\\[3pt]
False $m$, true $u$ & Supply a computation. Field $s$ is the predecessor,
$u$ the initializer, and $t$ the update.\\[3pt]
False $m$, false $u$ & Introduce the unary paired goal with witness $t$.
Canonical $u=0$ selects progress for a literal equation;
canonical $u=[]$ selects its zero-case premise. Field $s$ is unused.\\
\bottomrule
\end{tabularx}
\end{center}

Every canonical absent major is \texttt{[]}; it closes immediately
against the initial truth assumption. Scalar zero is deliberately not
an alternative absent major. The next section proves why this detail
makes the unguarded matcher safe.

\subsection{Cost and verification}
Relative to their corresponding frozen predecessors, the explicit and
byte profiles each remove 16 normalized bytes, one token, and one AST
node. The generic profile removes 20 bytes, one token, and one AST node.
The changes to the mathematical protocol, rather than unrelated import
or driver edits, account for these matched differences.

Against the independent old envelope, the new byte profile goes from
558 to 542 bytes, a reduction of approximately $2.87\%$. The explicit
list-based program is 550 bytes. The generic traversal reaches 261
tokens and 295 AST nodes, improving each old structural minimum by one;
its 578 bytes do not beat 558. Consequently no single simultaneous
three-coordinate record is asserted.

The normalizer removes actual docstrings, applies \texttt{ast.unparse},
and appends a newline. It counts the entire module. The delivered run
uses CPython 3.12.14 and zlib 1.3.2. Gzip measures
\texttt{gzip.compress(source, compresslevel=9, mtime=0)}, including its
wrapper. Raw and normalized measurements, source hashes, and the actual
gzip streams are supplied. Compression is an informational measurement.

The offline checker suite runs 46 positive, negative, alias, and malformed
cases through the three compact programs and the readable expansion,
under both ordinary Python and \texttt{-O}: 368 expected process results.
Two additional mutation controls replace the correct mode comparison
with the tempting wrong one, and demonstrate acceptance of a forbidden
compound-progress certificate. The corrected checker rejects that
certificate. This checks the actual restriction whose loss would matter,
rather than only a generic malformed input.

The finite arithmetic diagnostics separately evaluate the displayed
$B$ macros and two monus equations on inputs through eight, including
the intermediate equalities of every active gated-complement step.
They perform 1,750 assertions and distinguish seven failure mutants.
Those diagnostics are not a full proof checker for $B+d$ and are not
included in the compact trusted-source measurements.

% Independent audit fragment for the corrected array-truth interface.
% Reviewed files: kernel_notes/array_truth_checker{,_bytes,_joint}.py.
% IMPORTANT: the final mode comparison is with u, not s.
% Requires only ordinary LaTeX plus amsmath/amssymb.

\section{Why the absent major can be a proved empty array}
\label{sec:neutral-major-audit}

This section audits the new control-flow change on raw input, including
noncanonical JSON.  It uses the inherited sorted arity decoder and
semantic-opening lemmas for the numerical and logical rules.  It does not
replace those lemmas with finite testing or claim proof-assistant
certification.

\subsection{The exact change}

The initial context is the one-element list \texttt{[[]]}.  Its sole
formula, the empty array, denotes $0=0$ under the inherited decoder.
At every non-hypothesis command the checker now executes
\begin{verbatim}
    m, s, u, t = q()
    c(m, g)
\end{verbatim}
without testing whether \texttt{m} is nonempty.  A canonical command with
no substantive major puts \texttt{[]} in this field.  Its premise closes
immediately by hypothesis, without consuming a command or changing the
fresh-name floor.

The final choice of the retained assumption uses
\begin{verbatim}
    (m or len(a) % 2 == u) and k
\end{verbatim}
instead of comparing the arity parity with \texttt{m}.  In an
introduction command, \texttt{u=0} selects equation progress and
\texttt{u=[]} selects the zero-case premise.  This reuse is legitimate
because introduction is reached only when both \texttt{m} and
\texttt{u} are false-valued.  Computation still requires a truthy
initializer \texttt{u}; its missing literal-zero cases have the
inherited certified-zero repair.

The source is deliberately left with the old matcher
\begin{verbatim}
    m in ([s, u], not s and [t] * 3)
\end{verbatim}
and no added \texttt{m and} guard.  As an isolated Python expression this
would be unsafe: if \texttt{s} is truthy, its second alternative is
\texttt{False}, which compares equal to numeric zero.  The following
whole-program invariant proves that this dangerous major cannot reach
the matcher.

\subsection{Two facts about the root opener}

Call a Python value \emph{numeric} here when it is an integer, a float,
or a Boolean returned by the JSON loader.  The finite Python execution
also admits the loader's optional nonfinite float values.  No custom
numeric class is supplied by a certificate.

Write $O_k(r;t)$ for the opener on the numeric value $r$, with image $t$
and its default formula-root depth zero.  Direct evaluation of the
program gives
\[
 O_k(r;t)=
 \begin{cases}
 r,&k\le r<0,\\
 0,&\text{otherwise}.
 \end{cases}
 \tag{Root scalar}
\]
In particular this result is numeric and independent of the image $t$.
The formula covers negative allocated names and negative fractional
names without identifying different numerical values.  Signed zero and
Booleans have Python's ordinary numerical equality behavior.  A NaN
fails the interval test and returns zero; the relevant infinite values
also return zero.  None of these results is a list.

A second fact is that a successful root opening of a list returns a
list.  Under either strict list predicate, no other input type can
return a numeric result except a numeric input.  Under the generic
iterable opener, a successful iterable traversal also returns a list;
if traversal fails, the numerical comparison cannot succeed on a list,
string, dictionary, or null.  Thus in every profile:
\begin{quote}
If a root opening returns a numeric result, its argument was numeric;
if it succeeds on a list, its result is a list.
\end{quote}
These statements concern finite successful evaluations.  They assert
neither arbitrary revalidation idempotence at ignored annotations nor
unlimited-stack termination for the generic iterable profile.

\subsection{No numeric goal can be proved from list-valued assumptions}

\paragraph{Lemma.}
For any initial integer fresh-name floor and any finite remaining
certificate, there is no finite successful invocation
\texttt{c(a,g)} whose raw goal \texttt{a} is numeric and whose every
context element is a list.

\emph{Proof.}
Suppose otherwise, and select such a successful invocation with the
fewest calls to \texttt{c} in its completed recursive trace.  Its first
hypothesis check cannot close: the root opening of its numeric goal is
numeric, and no numerical JSON value equals a Python list.

The invocation must therefore consume a command and successfully return
from \texttt{c(m,g)}.  If \texttt{m} were numeric, this would be a
strictly shorter successful invocation of the forbidden kind.  Thus the
major is not numeric.

In either strict profile, a successfully checked nonnumeric major is a
list.  In the generic profile an empty string or an empty dictionary can
also open to the empty list.  A nonempty dictionary can occasionally
have its iterable alias proved, but it cannot occur as a truthy major
in the selected completed parent invocation: the subsequent expression
\texttt{m * 3} raises a type error.  A nonempty string has no finite
successful opening.  Consequently every truthy major that survives to
the continuation exports a list, and every surviving false-valued
major is an empty list or one of those harmless empty iterable aliases.

Consider the possible change of goal.  If the matcher does not fire,
the goal stays numeric.  If it fires, its target test requires
\texttt{i(t,u) == a}.  The second root-opening fact therefore implies
that \texttt{t} is numeric.  By (Root scalar), \texttt{i(t,s)} and
\texttt{i(t,u)} are the same root result.  Both the Boolean-expression
and the explicit-conditional implementations consequently retain a
numeric goal.  The Boolean implementation may retain the original raw
zero alias when the transported result is false-valued; that makes no
difference to this assertion.

If the major is truthy, the generated hypothesis is the list
\texttt{m * 3}.  If the major is false-valued and \texttt{u} is truthy,
the computation expression constructs a two-element list.  A successful
opening of either hypothesis is a list.  The final recursive call then
has a numeric root-opened goal and an entirely list-valued context.  It
is a proper subtrace of the selected invocation, contradicting its
minimality.

The only remaining case is introduction.  Its expression
\texttt{i(*a,t,2)} must unpack the numeric raw goal as an iterable,
which raises a type error.  That case cannot belong to a successful
trace.  All possibilities have been excluded.\hfill$\square$

The minimization ranges over all starting floors and stacks.  It is
therefore unaffected by a recursively checked major decreasing the
global floor before its parent resumes.  The argument uses finite
completed traces, not a presumption that recursive checking always
terminates.

\subsection{Every reached command has a list-valued context}

Initially the only assumption is the list \texttt{[]}.  Suppose a
command is reached with list-valued assumptions and its major check
returns.  The lemma excludes a numeric major.  A truthy surviving major
exports a list; a computation constructs a list.  In introduction the
assignment is
\begin{verbatim}
    h = a = i(*a, t, 2)
\end{verbatim}
in mathematical order.  If its result is numeric, the later
\texttt{len(a)} raises a type error before the continuation context is
constructed, since this branch has false-valued \texttt{m}.  Such a
scalar matrix does \emph{not} insert a scalar-zero assumption.  If the
assignment yields a list, successful opening again yields a list.

Both recursive premises thus receive list-valued contexts whenever
execution reaches them.  The expression \texttt{g + [new]} constructs a
new list; a premise cannot mutate its caller's context.  Induction on
the reached calls proves the invariant.  The tuple-based explanatory
variant has the corresponding immutability property.

It follows that no numeric zero, Boolean \texttt{False}, floating zero,
truthy numeric scalar, or newly allocated numeric name can be a
successfully verified major at a reached command.  In particular a
closed scalar-zero goal is rejected by this interface.  The canonical
empty-array truth is accepted instead; mathematical completeness is
claimed for canonical encodings, not for every redundant truth alias.

\subsection{The unguarded matcher is safe on every reached continuation}

If \texttt{s} is false-valued, both matcher alternatives are nonempty
lists, so no false-valued major can match.  If \texttt{s} is truthy, the
second alternative is \texttt{False}.  A false-valued major reaching
this point is an empty list, or in the generic profile an empty string
or dictionary.  None equals \texttt{False}; a numeric major has already
been excluded.  Therefore, at every point reached after a successful
major check, the matcher has exactly the same truth value as
\begin{verbatim}
    m and m in ([s, u], not s and [t] * 3)
\end{verbatim}
even though this identity is false on unrestricted Python tuples of
values.  No sentinel and no extra source guard are needed.

The proof depends on the initial context being \texttt{[[]]}, not
\texttt{[0]}, and on premise contexts being freshly constructed.
Silently changing either detail would invalidate this argument.

\subsection{Why the relocated induction mode must be the initializer field}

The corrected selector compares arity parity with \texttt{u}.  In an
introduction branch this field is necessarily false-valued.  No
false-valued JSON value compares equal to one.  Thus a meaningful
compound matrix, of odd arity one or three, cannot receive a progress
assumption.  A literal equation has even arity two and receives
progress precisely for a zero-like mode such as \texttt{0}; an empty
array mode chooses the zero instance.  Surplus arities retain their
inherited constant interpretations.  Scalar matrices cannot survive
the shape check discussed above.

Comparing instead with the unused \texttt{s} field would be wrong.
That field is unrestricted in introduction, so choosing
\texttt{s=1} would enable progress for odd-arity compound matrices.
The distinction is semantic, not merely a preferred serialization.

In the computation branch \texttt{u} is truthy.  Its new comparison can
choose a fresh-name opening for some noncanonical scalar initializer
aliases.  This is handled by the inherited zero-instance discharge
lemma, not by claiming that every generated open hypothesis is
uniformly true.  If the generated hypothesis is $H(\vec b)$, its zero
completion $H(\vec0)$ is a computation theorem, while the old-floor
validation makes the finite family of newly admitted names $\vec b$
fresh for the retained conclusion and context.  Substituting zero for
those names and cutting in the computation theorem discharges the
extra hypothesis.  Negative fractional names in the newly admitted
interval are included separately in this finite substitution.

\subsection{Conservativity and canonical completeness}

The initial empty-array formula is already a theorem.  For example the
permitted nonempty-seed computation gives
\[
 C=0,\qquad C:=I(S(S0),S0;v.0).
\]
Use this equation as both the equality major and the premise of
transport through $F(w)\equiv(w=0)$ to obtain $0=0$.  Adding that fixed
theorem to a proof context is conservative by substitution of its
finite derivation for hypothesis uses.  No new object-level truth
predicate or logical axiom has been introduced.

For canonical certificates, substantive major commands retain their
old fields.  A computation uses the absent-major field \texttt{[]} and
its existing nonempty initializer.  An introduction uses absent major
\texttt{[]}, with \texttt{u=0} for equation progress and
\texttt{u=[]} for cases; its witness remains in \texttt{t}.
The added empty major closes without a command or name allocation, so
the remaining canonical proof simulation is unchanged.  First applying
the existing certified-zero repair supplies every unrestricted
computation instance.  This gives a finite certificate for every
canonical derivation of the inherited native calculus.

Conversely, the invariant makes each reached matcher a legitimate
guarded matcher.  The remaining commands use the inherited
transport, closure, paired-introduction/elimination, and zero-instance
computation arguments, with the corrected equation-only mode test.
Discharging the initial $0=0$ theorem yields a derivation in that same
calculus.  Its established whole-domain iPRA interpretations therefore
transfer without changing the numerical domain, equality, or object
language.  The new result concerns a proof-checker presentation, not a
new primitive-recursive arithmetic foundation.

\subsection{Raw JSON types and the generic profile}

The strict profiles recognize only Python lists as containers.  The
short predicate \texttt{a * 0 == []} has this effect for JSON values:
lists pass; strings and numbers give different zero products;
dictionaries and null raise.  Its numerical fallback rejects strings.
The root invariant covers all numbers and Booleans that reach it,
including nonfinite values accepted by the default Python loader.

The generic profile has the inherited additional aliases.  An empty
string and an empty dictionary open to \texttt{[]}; a dictionary with
only the empty-string key opens to \texttt{[[]]}, with its value ignored.
A truthy dictionary cannot survive export as a major because
\texttt{m * 3} raises.  Nonempty strings and nonempty dictionary keys
have no finite successful opening.  Actual CPython rejects them through
its recursion limit and subsequent type failure, while an ideal
unlimited-stack execution can diverge.  The raw soundness argument
above covers every finite successful run of this profile; it does not
assert ideal termination on every malformed finite JSON input.

Larger arrays remain arrays throughout opening, irrespective of their
ignored annotations or depth increments.  They therefore do not
threaten the scalar-versus-list invariant.  Their formula meanings and
freshness behavior remain subject to the inherited sorted-decoder
proof.  Physical time, memory, and stack limits can reject sufficiently
large genuine proofs in every profile; the completeness claim has the
same adequate-resource, canonical-proof interpretation as its
predecessors.


\section{Why zero-source transport alone is insufficient}
\label{sec:zero-obstruction}

The native checker already shares the substitution path between equality
transport and closure specialization from a zero instance. It is tempting
to simplify the remaining match by allowing transport only from an atom
$0=t$. In the original presentation with computation leaves and no
separately granted truth axiom, this loses even $0=0$.
The counterinterpretation below does not apply verbatim to the new
checker while retaining its initial truth context $[\,[]\,]$: that
additional premise already supplies $0=0$ and is false in the
counterinterpretation. Thus this obstruction does not rule out
zero-source transport together with a separately primitive reflexivity
or truth axiom.

\begin{proposition}
In the initialized-iterator native calculus, replacing Tr by
\[
 \frac{\Gamma\vdash0=t\quad\Gamma\vdash F(0)}
      {\Gamma\vdash F(t)}
\]
does not suffice to derive $0=0$ from the computation schemas and the
remaining native rules.
\end{proposition}
\begin{proof}
Give a proof-theoretic counterinterpretation on $\{0,1\}$. Interpret
$S0=S1=1$, and interpret every iterator description as the constant
function $1$. Interpret the native atom $s=t$ as the relation
\[
 R(s,t)\quad\Longleftrightarrow\quad s=1.
\]
This auxiliary relation is \emph{not actual equality}; we are proving
nonderivability in a weakened calculus, not displaying an ordinary
equality model of iPRA. Interpret $Q,D$ by their ordinary formula meanings.

Every computation atom is true, since its left side is an iterator and
therefore has value one. Every $0=t$ atom is false, so the restricted
transport rule is vacuously sound. Every domain element is zero or a
successor. Moreover induction is valid for every predicate on this
domain: its base gives the value at zero and its successor step at zero
gives the value at one. Thus DI-case and even unrestricted DI-progress
are sound in this auxiliary interpretation. Hyp, QE, DE, and structural
rules are sound as well. Induction on finite derivations preserves truth.
But $R(0,0)$ is false. Hence no derivation of $0=0$ exists.
\end{proof}

Unrestricted transport does derive it. Take the computation
$H:I(0,0;v.0)=0$ and the motive $F(w)\equiv w=0$. Use $H$ as both the
major equality and the proof of $F(I(0,0;v.0))$. Tr gives $F(0)$.
If the checker restricts literal-zero initializers, first obtain $H$
by the certified-zero repair. The issue is the transport restriction,
not that optional direct-computation restriction.

The accompanying finite program checks the truth-table portions of this
counterinterpretation. The mathematical argument above, rather than a
bounded search for proofs, establishes nonderivability. It rules out
this particular fusion; it is not a lower bound on all iPRA checkers or
on all possible recodings of equality.

\section{Further research questions}

\begin{enumerate}
\item \textbf{A measured checker for the new numerical basis.}
Can the tied-seed and monus computation schemas be checked more simply
than arbitrary initialized scalar iteration, once every logical rule,
scope check, parser, and driver is charged? The present small checker
does not answer this question: it intentionally retains the inherited
$I$ language.

\item \textbf{One program below all coordinates.}
Can a structurally terminating explicit-parser implementation beat the
new family envelope in a single source file? The generic and byte
profiles should remain separate observations. Gzip size should remain
informational rather than the objective.

\item \textbf{Mechanical verification of the bootstrap.}
The most useful immediate formalization target is the gated-complement
identity with its fixed parameter, followed by the native contextual-case
lemma using $H$. The bounded compiler's cap coherence and finite-case
induction preservation should then be verified, with actual source and
target syntax and both return maps.

\item \textbf{The minimum-only candidate.}
Replacing the equality primitive by three recursive minimum equations
was investigated. The term
$\operatorname{parity}(\min(2x,2y+1))$ has the intended comparator
meaning in the standard model. Obtaining all needed monus equations
and whole-domain reflection from that basis, without shifted-parameter
induction, remains open here. The positive theorem in this report uses
the explicitly stated monus equations instead.

\item \textbf{General neutral-premise transformations.}
The checker reduction suggests a general method: adjoin a derivable
neutral premise so a partial proof interface becomes total. Determine
which other native rule interfaces admit this transformation without
adding parser complexity or allowing malformed values to imitate a
genuine inference. The numeric-major exclusion and the false-valued
mode field are substantive obligations in the present example.

\item \textbf{Proof length and readability.}
The small trusted source should also be compared with certificate size,
maximum call depth, compilation cost, and the complexity of its soundness
proof. The neutral-premise transformation adds a proof call that closes
immediately, but does not require a new serialized subproof for its
neutral major. Triplicated closure matrices remain a separate source
of certificate growth.
\end{enumerate}

\section{Reproduction and trust boundary}
Run \texttt{python3 verify.py} from the extracted directory to reproduce
the source measurements, positive and negative certificates in ordinary
and optimized Python, finite arithmetic diagnostics, and transport
counterinterpretation. Run \texttt{python3 build.py} to compile this
report. The recorded runtime and results are in
\file{evidence/verification.json}.

The complete theorem checkers import only the standard JSON library.
Producers, the readable expansion, diagnostics, measurement scripts, and
build tools are outside their trusted source. No input Python expression
is evaluated. Acceptance requires normal termination and exactly
\texttt{ACCEPT} on stdout; an exception or nonzero exit is rejection.
Unused certificate prefixes are inert, matching the incumbent interface.

Finite tests support implementation fidelity; the written invariants,
derivation translations, and interpretation proofs carry the general
claims. Physical stack and memory limits restrict every executable
implementation; canonical proof completeness concerns its
unbounded-resource mathematical abstraction. The generic traversal has
the additional malformed-string qualification already disclosed.

\clearpage\appendix
\section{Attributed bounded initialization and finite function graphs}
\label{app:bounded}
This appendix reproduces the bounded-initializer draft's construction
\cite{bounded}, beginning after its equality characteristic has become
available. It is included so the interpretation claim is inspectable
without retrieving another archive. The inherited monus term is written
$\widehat d$ to distinguish it from our primitive $d$; numerical constant-body
recursion is $C$ rather than $D$. The local hit function $H$ below is unrelated
to the bounded selector called $H$ in the native bootstrap.
All other mathematical content is retained, with notation and cross-reference
repairs. The hypotheses about $\epsilon$ are supplied by our proved
reflection theorem, rather than new axioms of $T_{Bd}$.

\subsection{Order and bounded arithmetic}

For a Boolean $b$, let $J(a,b):=B(Zb;u,v.a)$, so
$J(a,0)=a$ and $J(a,1)=Sa$.  Define
\begin{align*}
 H(n,y)&:=L(n;r,v.J(v,\epsilon(r,y))),\\
 \widehat d(x,y)&:=L(x;u,v.J(v,H(Su,y))),\\
 \operatorname{mn}(x,y)&:=\widehat d(x,\widehat d(x,y)),\\
 \operatorname{db}(n)&:=L(n;u,v.S(Sv)),\\
 \operatorname{ad}_b(a,c)&:=
 \widehat d\bigl(\widehat d(\operatorname{db}(b),\widehat d(b,a)),\widehat d(b,c)\bigr).
\end{align*}
All five are terms in the original finite-description language.
Use $x\leq y$ as the quantifier-free abbreviation $\widehat d(x,y)=0$.
The elementary order lemma proves that this is a discrete linear order,
that $\widehat d$ is its truncated difference, and that $\operatorname{mn}$
is its minimum.  It also proves, entirely by quantifier-free induction,
the following bounded-addition identities:
\begin{align}
 a\leq b&\ \Longrightarrow\ \operatorname{ad}_b(a,0)=a,
 \label{eq:ib-add0}\\
 a,Sc\leq b&\ \Longrightarrow\
 \operatorname{ad}_b(a,Sc)=S\operatorname{ad}_b(a,c),
 \label{eq:ib-addS}\\
 a,c\leq b&\ \Longrightarrow\
 \operatorname{ad}_b(a,c)\leq\operatorname{db}(b),
 \label{eq:ib-addbound}\\
 \operatorname{ad}_b(b,b)&=\operatorname{db}(b).
 \label{eq:ib-adddiag}
\end{align}
The function $c\mapsto\operatorname{ad}_b(a,c)$ is nondecreasing on
$c\leq b$, when $a\leq b$.  Only these bounded identities are used;
there is no uncharged total binary addition term in the source language.
The detailed inductions are given in Appendix~\ref{sec:internal-order-foundations}.

Define
\[
 \operatorname{sel}_b(c,a,e):=
 \operatorname{ad}_b\bigl(G(Zc,a),G(c,e)\bigr).
\]
If $a,e\leq b$, coverage and the gate equations show that this returns
$a$ at $c=0$ and $e$ at every successor control.  At most one summand
is nonzero, so the result is again at most $b$.

For any fixed earlier update $h$, define the clipped initialized recursor
\[
 R_b(n,a;u,v.h):=
 L\bigl(Sn;r,s.
 \operatorname{sel}_b\bigl(r,\operatorname{mn}(a,b),
        \operatorname{mn}(h[Pr/u,s/v],b)\bigr)\bigr).
 \tag{Clipped recursion}\label{eq:ib-clipped-recursion}
\]
The gate, minimum, and $L$ equations give
\begin{align}
 R_b(0,a;h)&=\operatorname{mn}(a,b),\label{eq:ib-R0}\\
 R_b(Sn,a;h)&=
 \operatorname{mn}\bigl(h(n,R_b(n,a;h)),b\bigr).
 \label{eq:ib-RS}
\end{align}
Quantifier-free induction gives $R_b(n,a;h)\leq b$ for every counter.
The cap $b$ is a passive parameter throughout the recursion.
This supplies arbitrary initialization whenever an independent cap
has first been fixed.  It is compatible with the one-coordinate growth
obstruction: the output is bounded in the cap coordinate.

\section{A compiler with internal bounds and cap coherence}
\label{sec:internal-compiler}

Fix the usual finite primitive-recursive descriptions built from zero,
successor, projections, composition, and primitive recursion.  The
following theorem is an external construction on descriptions, producing
finite source terms and finite source proofs.  It does not add a uniform
evaluation function for arbitrary description codes.

\begin{theorem}[Internal bounded realization]
\label{thm:internal-bounded-realization}
For every fixed $k$-ary primitive-recursive description $f$, one can
construct source terms
\[
                     \tau_f(b,x_1,\ldots,x_k),\qquad M_f(b)
\]
and proofs in $T_{B\epsilon}$ with the following properties.
\begin{enumerate}
\item $M_f$ is nondecreasing and $b\leq M_f(b)$.
\item The bound $\tau_f(b,\vec x)\leq M_f(b)$ holds for all inputs,
      whether or not the inputs are bounded by $b$.
\item If every $x_i\leq b$ and every $x_i\leq c$, then
      $\tau_f(b,\vec x)=\tau_f(c,\vec x)$.
\item The finite-case graphs built from these terms satisfy the defining
      equations of the description $f$, as proved below.
\end{enumerate}
All comparisons are the quantifier-free source abbreviations just
defined.  The construction produces proofs on the whole domain,
not just correctness on standard numerals.
\end{theorem}

\begin{proof}
We construct the first three items by induction on the fixed description.
The fourth will then follow from the actual defining formulas.

\paragraph{Initial descriptions.}
For zero use $\tau_0=0$; for a projection use
$\tau_{\pi_i}(b,\vec x)=\operatorname{mn}(x_i,b)$; for successor use
$\tau_S(b,x)=S\operatorname{mn}(x,b)$.  The common bound $M(b)=Sb$
works.  The minimum identities prove bounds and coherence immediately.

\paragraph{Composition.}
Suppose $f(\vec x)=h(g_1(\vec x),\ldots,g_m(\vec x))$.
Write $G_i=M_{g_i}$ and $H=M_h$.  Put
\[
 Q=G_1\circ\cdots\circ G_m,
 \qquad M_f=H\circ Q,
\]
with $Q$ the identity if $m=0$.  Since all $G_i$ are nondecreasing
and inflationary, $G_i(b)\leq Q(b)$ for each $i$.
Define
\[
 \tau_f(b,\vec x)=
 \tau_h\bigl(Q(b),\tau_{g_1}(b,\vec x),\ldots,
                         \tau_{g_m}(b,\vec x)\bigr).
\]
The global bound follows from the bound for $\tau_h$.  For coherence,
the inner values agree at any two input caps by their induction
hypotheses.  They are bounded by both $Q(b)$ and $Q(c)$; applying
coherence for $h$ proves the result.  Monotonicity and inflationarity
of $M_f$ follow by composition.

\paragraph{Primitive recursion and its majorant.}
Suppose
\[
 f(0,\vec x)=g(\vec x),\qquad
 f(Sn,\vec x)=h(n,f(n,\vec x),\vec x).
\]
Let $G=M_g$, $H=M_h$, and put
\[
                      F=H\circ G\circ S.
\]
This is a source unary term with internal proofs of monotonicity and
$Sz\leq F(z)$.  Let
\[
 I_F(t):=L(t;u,v.F(v)),\qquad
 C_b:=I_F(S\operatorname{db}(b)),\qquad M_f(b):=C_b.
 \tag{Majorant}\label{eq:ib-majorant}
\]
Equation induction gives $I_F(0)=0$, $I_F(St)=F(I_F(t))$,
and $t\leq I_F(t)$.  Its monotonicity is a quantifier-free induction
using order discreteness: if $s\leq St$, either $s=St$ or $s\leq t$.
Consequently $b\mapsto C_b$ is nondecreasing and inflationary.
This construction uses only successor, composition, and zero-start
iteration for its majorants.

For $j\leq b$, use the following actual source term for the stage bound:
\[
                 K_b(j):=I_F\bigl(S\operatorname{ad}_b(b,j)\bigr).
 \tag{Stage bound}\label{eq:ib-stage-bound}
\]
The bounded-addition identities prove
\begin{align}
 b,G(b)&\leq K_b(0),\label{eq:ib-Kbase}\\
 Sj\leq b&\ \Longrightarrow\ K_b(Sj)=F(K_b(j)),
 \label{eq:ib-Kstep}\\
 j\leq b&\ \Longrightarrow\ b\leq K_b(j)\leq C_b.
 \label{eq:ib-Kcap}
\end{align}
For example, $K_b(0)=I_F(Sb)=F(I_F(b))\geq G(b)$ because
$I_F(b)\geq b$, $G$ is nondecreasing, and $H$ is inflationary.
The upper bound uses $\operatorname{ad}_b(b,j)\leq\operatorname{db}(b)$.
Also $H(z)\leq F(z)$, by inflationarity of $G$ and monotonicity of $H$.

Write $n_b=\operatorname{mn}(n,b)$ and
$\vec x_b=(\operatorname{mn}(x_i,b))_i$.  Define
\[
 \tau_f(b,n,\vec x):=
 R_{C_b}\bigl(n_b,\tau_g(b,\vec x_b);
       u,v.\tau_h(C_b,u,v,\vec x_b)\bigr).
 \tag{Compiled recursion}\label{eq:ib-compiled-recursion}
\]
Its global bound is immediate from the clipped-recursion bound.

\paragraph{No clipping on bounded inputs.}
Fix passive parameters $\vec x\leq b$ and let
\[
 r_b(j):=R_{C_b}\bigl(j,\tau_g(b,\vec x);
                u,v.\tau_h(C_b,u,v,\vec x)\bigr).
\]
We prove by quantifier-free induction on $j$ that
\[
                       j\leq b\ \longrightarrow\ r_b(j)\leq K_b(j).
 \tag{Invariant}\label{eq:ib-growth-invariant}
\]
At zero, $\tau_g(b,\vec x)\leq G(b)\leq K_b(0)\leq C_b$,
so the seed is not clipped.  For the step assume $Sj\leq b$.
Then the induction hypothesis applies at $j$; all arguments
$j,r_b(j),\vec x$ of the step function are bounded by $K_b(j)$,
as well as by $C_b$.  Coherence for the already constructed description
$h$ gives
\[
 \tau_h(C_b,j,r_b(j),\vec x)
 =\tau_h(K_b(j),j,r_b(j),\vec x)
 \leq H(K_b(j))\leq F(K_b(j))=K_b(Sj)\leq C_b.
\]
Thus its value is not clipped and the invariant follows.  The same
argument establishes the actual untruncated step identity
\[
 Sj\leq b\ \Longrightarrow\
 r_b(Sj)=\tau_h(C_b,j,r_b(j),\vec x).
 \tag{Unclipped step}\label{eq:ib-unclipped-step}
\]
The matrix of the induction is quantifier free, and the input cap and
all passive parameters remain fixed throughout it.

\paragraph{Coherence between two arbitrary caps.}
Fix $b,c$ and inputs $n,\vec x$ bounded by both.  After the minimum
identities remove their input clipping, compare the two runs $r_b,r_c$.
Quantifier-free induction on $j$ proves
\[
             (j\leq b\wedge j\leq c)\ \longrightarrow\ r_b(j)=r_c(j).
 \tag{Two-cap invariant}\label{eq:ib-two-cap-invariant}
\]
The base uses coherence for $g$ and the already proved absence of seed
clipping.  If $Sj\leq b,c$, the two states at $j$ are equal by the
induction hypothesis and bounded by both $C_b,C_c$.  The other step
arguments are bounded by both caps too.  Coherence for $h$ equates
the two raw step values.  Formula~\eqref{eq:ib-unclipped-step}
removes their clipping, giving the successor equality.  At $j=n$
this is exactly the required coherence of $\tau_f$.
No induction on a universally quantified formula has occurred:
$b,c,n,\vec x$ are fixed parameters, and both displayed invariant
matrices are quantifier free.

Every recursive step body in this construction is the already compiled
original-language term $\tau_h$, inlined before it is inserted into $L$
or $R_b$. No recursion schema is extended for free to an opaque function
adjoined later by a graph definition. This distinction prevents a
conservative-definition argument from hiding a stronger recursion principle.

This completes the construction of the terms and internal proofs.
\end{proof}

\section{Function graphs and preservation of induction}
\label{sec:internal-function-graphs}

For a nonempty tuple $\vec x=(x_1,\ldots,x_k)$ let
\[
 C_i(\vec x):=
 \bigwedge_{j<i}x_j<x_i\ \wedge\
 \bigwedge_{j>i}x_j\leq x_i.
\]
This says that coordinate $i$ is the first largest one.  Discrete
linearity proves that these fixed finitely many quantifier-free cases
are exhaustive and disjoint.  Define
\[
 \Phi_f(\vec x,z):=
 \bigvee_{i=1}^{k}
 \bigl(C_i(\vec x)\wedge z=\tau_f(x_i,\vec x)\bigr).
 \tag{Graph}\label{eq:ib-graph}
\]
For a nullary description use $z=\tau_f(0)$.

\paragraph{A concrete example: binary addition.}
The bounded-addition term already illustrates the method without the full compiler. Its total single-valued graph is
\[
 \Phi_+(x,y,z):=
 \bigl(y\leq x\wedge z=\operatorname{ad}_x(x,y)\bigr)
 \ \vee\
 \bigl(x<y\wedge z=\operatorname{ad}_y(x,y)\bigr).
\]
On either case the selected cap bounds both arguments. The bounded-addition identities and coherence under changing a valid cap give the usual addition equations. Each possible output term is bounded by twice one input, while addition itself is not bounded by any unary function of one input. Thus graph interpretation succeeds exactly where a single-term definition is obstructed.

\begin{lemma}[Totality, uniqueness, and computation]
\label{lem:internal-graph-computation}
The theory proves that every $\Phi_f$ is total and single-valued.
In the definitional extension by its uniquely specified function $f^*$,
the defining composition and primitive-recursion equations of the
description $f$ hold.  The extension admits quantifier-free induction
for formulas containing finitely many of these defined functions.
\end{lemma}

\begin{proof}
Exhaustiveness provides a displayed value term on each case, proving
existence.  Disjointness proves uniqueness.  Moreover, for every
source cap $b$ bounding the arguments, coherence proves
\[
                              f^*(\vec x)=\tau_f(b,\vec x).
 \tag{Cap realization}\label{eq:ib-cap-realization}
\]
One proves this by the finite graph cases: on case $i$, both $b$ and
$x_i$ bound all arguments, so the corresponding values are equal.

Initial equations follow from the compiled initial descriptions.
For composition, use any cap $b$ bounding the common input tuple.
The compiled inner values are bounded by $Q(b)$, so applying
\eqref{eq:ib-cap-realization} at $Q(b)$ to the outer function proves
the composition equation directly from its compiler definition.

For recursion at counter zero, choose a graph case whose cap bounds
$0,\vec x$.  The seed bound prevents clipping, so the value is
$\tau_g(b,\vec x)=g^*(\vec x)$.  At a successor counter, choose a
case whose cap $b$ bounds $Sn,\vec x$.  Then $n\leq b$,
the unclipped-step identity applies, and
\eqref{eq:ib-cap-realization} identifies the previous state with
$f^*(n,\vec x)$.  All step arguments are bounded by $C_b$; the same
cap realization for $h$ identifies the update with
$h^*(n,f^*(n,\vec x),\vec x)$.  This proves the recursion equations.

It remains to justify the induction claim, rather than merely invoke
conservativity of definitional extensions.  Every occurrence of a
defined function has a graph given by finitely many quantifier-free
cases and explicit source value terms.  By external induction on a
finite term, it therefore has a finite family of pairs
$(A_j,t_j)$, with source quantifier-free guards $A_j$ and source
terms $t_j$, such that the guards cover all inputs and the term equals
$t_j$ on guard $A_j$.  At a composition of terms take the finite
Cartesian product of the argument cases and then apply the finitely
many cases of the outer function.  This introduces no object-language
existential quantifier.  Guards need not be disjoint after expansion;
their associated values agree on every overlap.

For a quantifier-free formula involving defined functions, take a
common finite refinement of the cases of all its terms.  On each
guard replace its terms by their source values and retain its actual
intuitionistic connectives.  The finite disjunction of these guarded
source formulas is quantifier free and is provably equivalent to the
original formula.  Covering cases justify the equivalence also for
implication and negation.  Applying source quantifier-free induction
to this equivalent formula transfers induction to the expanded
formula.  This is the essential finite-case step: an arbitrary
existential graph translation would not suffice.
\end{proof}

\section{Whole-domain bi-interpretability}
\label{sec:internal-biinterpretability}

\begin{theorem}[Arithmetic equivalence]
\label{thm:internal-biinterpretability}
The explicitly axiomatized theory $T_{B\epsilon}$ and full first-order
intuitionistic primitive recursive arithmetic are bi-interpretable.
Both interpretations use the whole numerical domain and actual
equality.  Their composite definable isomorphisms are identity maps.
In particular the theories are equiconsistent, and neither an appeal
to standard-model agreement nor induction on quantified matrices is
needed for this conclusion.
\end{theorem}

\begin{proof}
\paragraph{Interpretation in iPRA.}
Interpret $0,S$ ordinarily and $\epsilon$ as iPRA's primitive-recursive
equality characteristic.  Translate every fixed staged $B_h$ symbol
by primitive recursion with initial value $S h(0,0,\vec p)$ and update
$h(n,v,\vec p)$, first translating its earlier-stage body.  The two
computation equations become defining primitive-recursion equations.
The equality-test, separation, and contextual equation-induction
principles hold in iPRA.  Logical connectives, quantifiers, and genuine
falsity retain their ordinary intuitionistic meanings.

\paragraph{Interpretation of iPRA in the source.}
Interpret each fixed primitive-recursive function description by its
total unique graph $\Phi_f$ and use ordinary relational interpretation
of function terms, or equivalently work in the conservative
definitional extension by the functions $f^*$.  The graph computation
lemma proves every defining equation.  The source predecessor and
separation prove the successor boundary laws.  The finite-case
quantifier-free elimination in that lemma proves every translated
iPRA induction instance.  Thus all axioms and inference rules of
full iPRA are interpreted, with their intuitionistic logic unchanged.

\paragraph{The composite on the iPRA side.}
An externally fixed primitive-recursive description is uniquely
determined by its initial and step equations.  External induction on
its finite description, with ordinary quantifier-free induction on
the recursion counter at each recursion node, proves in iPRA that
its composite interpretation equals the original function.  The
necessary graph functions can be treated in their conservative
definitional extensions; finite-case elimination supplies any needed
induction instances.  Hence the identity on the domain preserves
every primitive-recursive function and actual equality.

\paragraph{The composite on the source side.}
Proceed by external induction on the finite stage of a source
description.  The ordinary $0,S$ operations are preserved.
The interpreted equality characteristic and the original $\epsilon$
are both Boolean characteristics of actual equality; their range and
reflection properties prove that they agree by the two Boolean cases.
For a fixed $B_h$, the earlier-stage induction hypothesis identifies
the interpreted update with its original source body $h$.  Both the
original function and its composite satisfy
\eqref{basis:b0} and \eqref{basis:bs} with that same body.  Induction
on the equation asserting their equality proves recurrence uniqueness
and identifies them at every counter and every passive parameter.
If the composite is expressed by defined graph functions, the
finite-case elimination lemma again legitimizes this induction in
the original theory.  Thus the identity map preserves every source
function symbol too.

Since all numerical operations and equality are preserved, external
induction on arbitrary formula syntax proves preservation of all
logical formulas under both composites.  Quantifiers range over the
same whole domain.  This establishes the two identity definable
isomorphisms and the claimed full bi-interpretation.  Genuine falsity
is preserved in both directions, giving equiconsistency.
\end{proof}


% Independent internal foundation fragment, 2026-11-10.
% Requires amsmath, amssymb, amsthm, and theorem/lemma environments.
% The surrounding development supplies P,Z,G,L,J, the numerical equality
% test delta, their endpoint laws, and derived quantifier-free induction.
% No primitive unrestricted addition, order, or constructor-coverage axiom
% is assumed in this fragment.

\section{Internal order, difference, and addition under a bound}
\label{sec:internal-order-foundations}

This attributed appendix supplies the arithmetic facts needed by the bounded compiler~\cite{bounded}. The notation $\widehat d$ distinguishes its constructed difference from the primitive $d$ of our new basis.
All its inductions are on quantifier-free formulas with the other displayed
variables held fixed as parameters.  In particular, an induction below is
never shorthand for induction on a universally quantified formula.  Use $\epsilon$ for the numerical equality test, with values zero and one, and use
the already derived terms and equations
\[
 P0=0,\quad P(Sx)=x,\qquad Z0=1,\quad Z(Sx)=0,
 \qquad G(0,a)=0,\quad G(Sx,a)=a,
\]
\[
 L_h(0)=0,\quad L_h(Sn)=h(n,L_h(n)),\qquad
 J(a,0)=a,\quad J(a,1)=Sa.
\]
Here $1=S0$.  Applying $P$ to an equality of successors gives successor
injectivity.  Applying $Z$ to $0=Sx$ contradicts $0\ne1$, so zero is
distinct from every successor.  The proved Boolean range and reflection laws consequently give
\[
 \epsilon(x,x)=1,\quad \epsilon(x,y)=\epsilon(y,x),\quad
 \epsilon(Sx,Sy)=\epsilon(x,y),\quad \epsilon(0,Sx)=0.
\]
Each identity follows by the two possible Boolean values and the reflection
law for $\epsilon$; these are derived identities, not extra axioms.

\begin{lemma}[Derived constructor coverage]
\label{lem:internal-coverage}
The theory proves $x=0\lor x=S(Px)$.
\end{lemma}
\begin{proof}
Apply the already derived quantifier-free induction to
$x=0\lor x=S(Px)$.  The base is immediate.  At a successor the second
disjunct follows directly from $P(Sx)=x$, independently of the induction
hypothesis.  Thus constructor coverage need not be a primitive axiom.
\end{proof}

\subsection{A hit term defines a discrete linear order}

Define
\[
 H(n,y):=L(n;r,v.J(v,\epsilon(r,y))).
\]
Its defining equations are
\[
 H(0,y)=0,\qquad H(Sn,y)=J(H(n,y),\epsilon(n,y)).
 \tag{H}\label{eq:internal-hit-recursion}
\]

\begin{lemma}[Hit identities and duality]
\label{lem:internal-hit}
The following identities and range assertion hold internally:
\begin{align*}
 H(Sn,0)&=1,& H(Sn,Sy)&=H(n,y),& H(n,n)&=0,\\
 H(Sn,n)&=1,& H(n,y)&=0\lor H(n,y)=1,\\
 H(Sy,x)&=Z(H(x,y)).
\end{align*}
Moreover a hit persists: $H(n,y)=1$ implies $H(Sn,y)=1$.
\end{lemma}
\begin{proof}
Equation induction on $n$ proves $H(Sn,0)=1$: its first step hits zero,
and every later equality test against zero is zero.  Equation induction
on $n$, with $y$ fixed, proves $H(Sn,Sy)=H(n,y)$, using
$\epsilon(Sn,Sy)=\epsilon(n,y)$.  A further equation induction gives
$H(n,n)=0$; the recurrence then gives $H(Sn,n)=1$.

For Boolean range use quantifier-free induction on $n$, with $y$ fixed.
If $\epsilon(n,y)=0$, the state is unchanged.  If $\epsilon(n,y)=1$, reflection
gives $n=y$, so the old state is $H(n,n)=0$ and the new state is one.
If $H(n,y)=1$, diagonal zero shows $n\ne y$, so the next equality test
vanishes.  This proves persistence.  The same case split shows
\[
 H(Sn,y)=1\ \Longrightarrow\ H(n,y)=1\ \lor\ n=y.
 \tag{New hit}\label{eq:internal-new-hit}
\]

Finally prove duality by equation induction on $x$, with $y$ fixed.
At zero both sides are one.  In the successor step the new left side is
$H(Sy,Sx)=H(y,x)$, whereas the old left side is
$J(H(y,x),\epsilon(y,x))$.  If $\epsilon(x,y)=0$, both relevant recurrences
leave their states unchanged, and the induction hypothesis gives the
required equality.  If $\epsilon(x,y)=1$, put $x=y$; the new left side is
zero and the new right side is $Z(H(Sx,x))=Z1=0$.
\end{proof}

Put
\[
 x<y\quad:\!\Longleftrightarrow\quad H(y,x)=1,
 \qquad
 x\le y\quad:\!\Longleftrightarrow\quad H(Sy,x)=1.
 \tag{Order}\label{eq:internal-order}
\]
These are quantifier-free abbreviations.

\begin{lemma}[Discrete linear-order laws]
\label{lem:internal-linear-order}
The relations in \eqref{eq:internal-order} form a discrete linear order
with least element zero.  In particular,
\[
 x\le y\ \Longleftrightarrow\ (x=y\lor x<y),\qquad
 x<Sy\ \Longleftrightarrow\ x\le y,
\]
successor preserves and reflects both comparisons, and
$x\le Sy$ with $x\ne Sy$ implies $x\le y$.
\end{lemma}
\begin{proof}
Diagonal identities give reflexivity of $\le$, irreflexivity of $<$,
and $x<Sx$.  The recurrence for $H(Sy,x)$, split on $\epsilon(y,x)$,
gives the first displayed equivalence.  Duality and Boolean range give
trichotomy: if $x\ne y$, then $H(Sy,x)=H(y,x)$, and $H(x,y)$ is its
Boolean complement.  Thus exactly one of $x<y$, $x=y$, $y<x$ holds.

For strict transitivity retain $H(y,x)=1$ and induct on $z$ in the
quantifier-free formula
\[
                 H(z,y)=1\ \longrightarrow\ H(z,x)=1.
\]
The base has false antecedent.  At the successor, the new-hit implication
\eqref{eq:internal-new-hit} gives either $H(z,y)=1$, when the induction
hypothesis followed by persistence applies, or $z=y$, when the retained
$H(y,x)=1$ followed by persistence applies.  Transitivity of $\le$ and
the mixed comparison laws now follow by splitting equality from strict
comparison.  Trichotomy gives antisymmetry.

The identities $H(Sn,Sy)=H(n,y)$ give successor preservation and
reflection.  The equivalence $x<Sy\leftrightarrow x\le y$ is immediate
from the definitions, and the asserted discrete predecessor step follows
from it and the first equivalence.  Finally $H(Sx,0)=1$ gives $0\le x$;
if $x\le0$, trichotomy or constructor coverage gives $x=0$.
\end{proof}

\subsection{Difference is an internal term, not an imported operation}

Define
\[
 \widehat d(x,y):=L(x;u,v.J(v,H(Su,y))).
\]
Thus
\[
 \widehat d(0,y)=0,
 \qquad \widehat d(Sx,y)=J(\widehat d(x,y),H(Sx,y)).
 \tag{d}\label{eq:internal-difference-recursion}
\]

\begin{lemma}[Difference laws]
\label{lem:internal-difference}
The theory proves
\begin{align*}
 \widehat d(x,0)&=x,& \widehat d(Sx,Sy)&=\widehat d(x,y),&\widehat d(x,x)&=0,\\
 \widehat d(x,y)=0&\ \Longleftrightarrow\ x\le y,
 &\widehat d(x,Sy)&=P(\widehat d(x,y)),&\widehat d(x,y)&\le x.
\end{align*}
Difference is monotone in its first argument and antitone in its second.
Furthermore
\begin{align*}
 y\le x&\ \Longrightarrow\ \widehat d(Sx,y)=S(\widehat d(x,y)),\\
 0<y\le x&\ \Longrightarrow\ \widehat d(x,Py)=S(\widehat d(x,y)).
\end{align*}
\end{lemma}
\begin{proof}
Equation induction gives $\widehat d(x,0)=x$ because $H(Sx,0)=1$.
For simultaneous successors, induct on $x$ with $y$ fixed.  The base
uses $H(S0,Sy)=0$.  At the successor use
$H(S(Sx),Sy)=H(Sx,y)$ and the induction hypothesis in
\eqref{eq:internal-difference-recursion}.

Quantifier-free induction on $x$, with $y$ fixed, proves
$\widehat d(x,y)=0\leftrightarrow H(x,y)=0$.
At the successor, if $H(Sx,y)=1$, the new difference is a successor and
is nonzero.  If $H(Sx,y)=0$, persistence implies $H(x,y)=0$, so the old
difference is zero and remains zero.  Duality identifies $H(x,y)=0$
with $H(Sy,x)=1$, proving the order characterization and $\widehat d(x,x)=0$.

Now induct on $x$ for $\widehat d(x,Sy)=P(\widehat d(x,y))$, with $y$ fixed.  Its successor
left side is $\widehat d(x,y)$.  The right side is
$P(J(\widehat d(x,y),H(Sx,y)))$.  When the Boolean is one, predecessor cancels
the successor; when it is zero, the preceding zero characterization
gives $\widehat d(x,y)=0$.  Both cases give the desired equality.

Quantifier-free induction on $x$ proves $\widehat d(x,y)\le x$: its successor
state is either the old state or its successor.  For first-argument
monotonicity induct on $v$ in
\[
       u\le v\ \longrightarrow\ \widehat d(u,z)\le \widehat d(v,z),
\]
with $u,z$ fixed.  At $Sv$, either $u=Sv$ or discreteness gives $u\le v$;
in the latter case use the induction hypothesis and the fact that the
new state is at least its old state.  The base uses $u\le0\to u=0$.

For second-argument antitonicity retain $y\le z$ and induct on $x$ in
$\widehat d(x,z)\le \widehat d(x,y)$.  If $H(Sx,z)=0$, the left state stays fixed and
the right state cannot decrease.  If $H(Sx,z)=1$, then $z\le x$;
transitivity gives $y\le x$, so both states increment and successor
preserves their old comparison.

The first final implication is the defining recurrence, since
$y\le x$ means $H(Sx,y)=1$.  For the second, coverage gives
$y=S(Py)$ when $y>0$.  The inequalities imply $Py<x$, so
$\widehat d(x,Py)\ne0$.  Functional coverage at this positive difference and
$\widehat d(x,S(Py))=P(\widehat d(x,Py))$ yield
$\widehat d(x,Py)=S(\widehat d(x,y))$.
\end{proof}

Define
\[
                      \operatorname{mn}(x,y):=\widehat d(x,\widehat d(x,y)).
\]

\begin{lemma}[Minimum]
\label{lem:internal-minimum}
If $x\le y$, then $\operatorname{mn}(x,y)=x$; if $y\le x$, then $\operatorname{mn}(x,y)=y$.
Consequently $\operatorname{mn}$ is the minimum operation for the internal linear order.
\end{lemma}
\begin{proof}
The first claim uses $\widehat d(x,y)=0$ and $\widehat d(x,0)=x$.
For the second induct on $y$ in
$y\le x\to \widehat d(x,\widehat d(x,y))=y$, with $x$ fixed.
The base is $\widehat d(x,x)=0$.  If $Sy\le x$, then $y<x$ and
$0<\widehat d(x,y)\le x$.  Apply the final difference implication to obtain
\[
 \widehat d(x,\widehat d(x,Sy))
   =\widehat d(x,P(\widehat d(x,y)))
   =S(\widehat d(x,\widehat d(x,y)))=Sy.
\]
The last equality uses the induction hypothesis, which applies because
$y\le x$.  Linearity completes the interpretation of $\operatorname{mn}$.
\end{proof}

\subsection{Addition with an explicit common bound}

Define the unary doubling term
\[
             \operatorname{db}(n):=L(n;u,v.S(Sv)).
\]
Thus $\operatorname{db}(0)=0$ and $\operatorname{db}(Sn)=S(S(\operatorname{db}(n)))$.  The notation $2n$ in the compiler
may be read as this term $\operatorname{db}(n)$; it does not invoke binary addition.

\begin{lemma}[Doubling identities]
\label{lem:internal-doubling}
The term $\operatorname{db}$ is monotone, $n\le \operatorname{db}(n)$, and $\widehat d(\operatorname{db}(n),n)=n$.
\end{lemma}
\begin{proof}
Quantifier-free induction gives $n\le \operatorname{db}(n)$ directly from the two
successors in the defining equation.  Monotonicity follows by induction
on the larger argument, splitting equality at its successor from the
discrete predecessor case, just as for first-argument monotonicity of $d$.
Equation induction proves the remaining identity.  At the successor,
\[
 \widehat d(\operatorname{db}(Sn),Sn)
   =\widehat d(S(S(\operatorname{db}(n))),Sn)=\widehat d(S(\operatorname{db}(n)),n)
   =S(\widehat d(\operatorname{db}(n),n))=Sn,
\]
where the third equality uses $n\le \operatorname{db}(n)$.
\end{proof}

Put
\[
 K_b(a):=\widehat d(\operatorname{db}(b),\widehat d(b,a)),\qquad
 \operatorname{ad}_b(a,c):=\widehat d(K_b(a),\widehat d(b,c)).
 \tag{Bounded addition}\label{eq:internal-bounded-addition}
\]
Every expression here is a finite term of the original language.

\begin{lemma}[Bounded-addition laws]
\label{lem:internal-bounded-addition}
Internally, the following hold:
\begin{enumerate}
\item If $a\le b$, then $\operatorname{ad}_b(a,0)=a$.
\item If $a\le b$ and $Sc\le b$, then
      $\operatorname{ad}_b(a,Sc)=S(\operatorname{ad}_b(a,c))$.
\item If $c\le b$, then $\operatorname{ad}_b(0,c)=c$.
\item $\operatorname{ad}_b(b,b)=\operatorname{db}(b)$ and $\operatorname{ad}_b(a,c)\le \operatorname{db}(b)$ for all $a,c$.
\item $\operatorname{ad}_b$ is monotone in its second argument.
\item If $a,c\le b$ and $a,c\le b'$, then
      $\operatorname{ad}_b(a,c)=\operatorname{ad}_{b'}(a,c)$.
\end{enumerate}
Thus it supplies addition on inputs bounded by $b$, without asserting an
unrestricted binary addition term.
\end{lemma}
\begin{proof}
Since $\widehat d(b,a)\le b$, antitonicity and doubling give
\[
            b=\widehat d(\operatorname{db}(b),b)\le K_b(a)\le \operatorname{db}(b).
 \tag{K bounds}\label{eq:internal-K-bounds}
\]
Also $K_b(0)=b$.  If $Sa\le b$, put $q=\widehat d(b,a)$.
Then $0<q\le b\le \operatorname{db}(b)$ and $\widehat d(b,Sa)=Pq$, so
\[
                  K_b(Sa)=\widehat d(\operatorname{db}(b),Pq)=S(K_b(a)).
\]
Quantifier-free induction on $a$ now proves the first assertion under
$a\le b$.  The base is $\widehat d(b,b)=0$; at the successor use the last
identity and $b\le K_b(a)$ to increment $\widehat d(K_b(a),b)$.

For the second assertion put $q=\widehat d(b,c)$.  The premise $Sc\le b$ gives
$0<q\le b\le K_b(a)$.  Hence
\[
 \operatorname{ad}_b(a,Sc)
   =\widehat d(K_b(a),Pq)=S(\widehat d(K_b(a),q))=S(\operatorname{ad}_b(a,c)).
\]
The third assertion follows by quantifier-free induction on $c\le b$
using the first two.  For the fourth, substitute $a=c=b$ and use
$\widehat d(b,b)=0$ twice; the upper bound follows from
$\widehat d(K_b(a),\widehat d(b,c))\le K_b(a)\le \operatorname{db}(b)$.

If $c\le c'$, antitonicity gives $\widehat d(b,c')\le \widehat d(b,c)$; applying
antitonicity again to the outer difference gives
$\operatorname{ad}_b(a,c)\le \operatorname{ad}_b(a,c')$.  Finally retain $a\le b,b'$ and induct on $c$
in the quantifier-free implication
\[
 c\le \operatorname{mn}(b,b')\ \longrightarrow\ \operatorname{ad}_b(a,c)=\operatorname{ad}_{b'}(a,c).
\]
The base gives $a$ on each side, and whenever $Sc\le \operatorname{mn}(b,b')$ the
successor laws apply to both sides and preserve the induction hypothesis.
\end{proof}

These proofs establish the numerical prerequisites for bounded selection
and clipped recursion entirely within the proposed theory.  Each use of
a case split is justified by numerical equality decidability, Boolean
range, or the derived linear-order and coverage laws.  No standard-model
correctness assertion and no hidden unrestricted addition operation is
used as an internal axiom.

\begin{corollary}[Identification with the new primitive monus]
In $T_{Bd}$, $\widehat d(x,y)=d(x,y)$ for every $x,y$.
\end{corollary}
\begin{proof}
Keep $x$ fixed and apply literal-equation induction in $y$.
Both terms have value $x$ at zero. Their successor values are the
predecessors of their previous values, by the proved difference lemma
and the primitive monus equation. Congruence of $P$ proves the step.
All occurrences of $\epsilon$ in $\widehat d$ are inlined as the fixed
$B+d$ term from the main text. Thus this is a source-language equation
induction, without an appeal to a graph-function extension.
\end{proof}

\clearpage
\section{Complete measured source}
The primary implementation uses an explicit list test. The byte profile
changes only that test to the inherited exact-JSON predicate
\texttt{a*0 == []}. The generic profile has the runtime qualification
stated in the implementation section. Each is a complete standalone
program, including JSON input and the acceptance driver.
\subsection{Explicit list test}
\lstinputlisting[language=Python]{code/checker.py}
\subsection{Byte profile}
\lstinputlisting[language=Python]{code/checker_bytes.py}
\subsection{Generic traversal profile}
\lstinputlisting[language=Python]{code/checker_joint.py}
\clearpage
\begin{thebibliography}{9}
\bibitem{snapshot}
Vladimir Reshetnikov, \emph{ProveIt}, \file{docs/drafts}, commit
\file{76a2651601d17b34d46753f189ea026a21ce8d29}, inspected 10 November 2026.
\url{https://github.com/VladimirReshetnikov/ProveIt/tree/76a2651601d17b34d46753f189ea026a21ce8d29/docs/drafts}.

\bibitem{certified}
ProveIt draft, \file{ipra_certified_zero_2026-11-10.zip}, in
snapshot~\cite{snapshot}. Source of the current smallest complete
checkers, initializer-based dispatch, the certified-zero normalization,
the source normalizer, and attributed inherited arithmetic and logic.

\bibitem{bounded}
ProveIt draft, \file{ipra_bounded_initializer_2026-11-10.zip}, in
snapshot~\cite{snapshot}. Source of the tied-seed $B+\epsilon$ basis,
bounded compiler, cap coherence, finite quantifier-free value graphs,
internal order lemmas, and the open reflection-replacement question.

\bibitem{local}
ProveIt draft, \file{ipra_local_scopes_2026-11-10.zip}, in
snapshot~\cite{snapshot}. Tied-seed macros, finite-profile obstructions,
and prior lexical-scope experiments.

\bibitem{arity}
ProveIt draft, \file{ipra_arity_calculus_2026-11-10.zip}, in
snapshot~\cite{snapshot}. Unary paired formulas, triplicated closures,
uniform arity-based depth, and conservative surplus-array semantics.

\bibitem{successor}
ProveIt draft, \file{ipra_successor_seed_completion_2026-11-10.zip}, in
snapshot~\cite{snapshot}. The finite-range guarded-case construction
adapted in the native bootstrap, with the Boolean-first selector supplied
here by different arithmetic terms.

\bibitem{affine}
ProveIt draft, \file{ipra_affine_supply_2026-11-10.zip}, in
snapshot~\cite{snapshot}, \file{evidence/theory_review.md}, section 5.
Contains an uncompleted $B+$monus suggestion and cautions about circular
equality-reflection arguments.

\bibitem{salehi}
S. Salehi, ``Decidable formulas of intuitionistic primitive recursive
arithmetic,'' \emph{Reports on Mathematical Logic} 36 (2002), 55--61.
\url{https://rml.tcs.uj.edu.pl/rml-36/36-SALEHI.pdf}.
\end{thebibliography}

\end{document}
